% !TeX root = ../../Book5.tex
% 纲要标题：## 第2章　局部环、有限模与局部—整体原则
% 纲要位置：几何学计划纲要.md，第 5070 行。
% 正文状态：初稿；尚未正式审核。

% 纲要内容开始
% 
% > 本章目标与证明范围：局部化、模运算和有限性一起服务于函数芽及局部构造。完备化的证明移至独立章，不在这里把一个形式邻域定义扩充成未建立的忠实平坦性。
% 
% 纲要内容结束

\chapter{局部环、有限模与局部—整体原则}
\label{g5:ch:02}
\BookChapterNumber{g5:ch:02}

\begin{chapterabstract}
	一个点附近的函数由局部化给出，函数之间的线性关系则需要用模记录。本章建立环与模的局部化、有限呈示和正合性，证明 Nakayama 引理，使剩余域上的线性代数能够控制有限模的生成元。支集、关联素理想与准素分解进一步辨认模存在的位置、零因子及嵌入分支。最后证明模和同态的局部检测原则，并将有限投射模与有限呈示的局部自由模联系起来；纤维计算同时说明有限性与平坦性条件的作用。
\end{chapterabstract}

\begin{learninggoals}
	\begin{itemize}
		\item 从分式等价关系构造局部化，证明泛性质、素理想对应，并计算局部环及剩余域。
		\item 用生成元和关系矩阵描述有限呈示模，区分局部化、张量积与 Hom 对正合列的不同作用。
		\item 证明并使用 Nakayama 引理，由剩余向量空间寻找极小生成元，并判断生成元能否在邻域内延展。
		\item 计算支集、零化子和关联素理想，完成所需的准素分解，区分不可约分支与嵌入信息。
		\item 证明有限呈示局部自由与有限投射的等价条件，检验秩和纤维维数，并解释恒定纤维维数反例。
	\end{itemize}
\end{learninggoals}

在第一章的两条坐标轴 \(xy=0\) 上，把 \(x\) 变成单位就会迫使 \(y=0\)；在原点附近，却不能除以 \(x\) 或 \(y\)，两条轴都必须保留。同一个坐标环，因所允许的分母不同，可以描述一条分支的一般位置，也可以描述分支相遇的局部结构。局部化把这种变化变成统一的代数运算。

\ProseParagraphBreak

还有一种信息不能只用函数环本身说明。考虑 \(k[x,y]\) 中由 \(x,y\) 生成的理想：它是函数的集合，又可以视为由两个元素生成的模，关系 \(yx-xy=0\) 则说明这两个生成元不是一组基。在原点之外，至少有一个生成元成为单位，这个理想局部等于整个函数环；在原点，它仍需要两个生成元。要理解这类随点变化的线性数据，需要同时追踪生成元、关系和局部化。

\ProseParagraphBreak

本章从交换含幺环、理想、商环和有限维线性代数出发，引入所需的模、张量积与正合列。第一章的 Hilbert 基定理和函数芽模型将分别提供 Noether 环的主要例子与局部化的几何解释。一般论证不要求基域代数闭，也会逐项区分任意模、有限生成模及有限呈示模；这些区别决定能否选择共同分母，能否在开邻域内同时控制有限多条关系。主要代数构造可结合\cite[第2、4节]{Matsumura1989}阅读，支集与关联素理想的具体阅读位置列在相应节中。

% !TeX root = ../../Book5.tex
% 纲要标题：### 2.1 局部化与局部环
% 纲要位置：几何学计划纲要.md，第 5074 行。
% 正文状态：初稿；尚未正式审核。

% 纲要内容开始
% 
% * 乘法集、分式及泛性质
% * 素理想处的局部化
% * 单位、剩余域与局部同态
% * 理想扩缩和饱和
% 
% 纲要内容结束

\section{局部化与局部环}
\label{g5:sec:0201}

在\cref{g5:prop:model-local-ring}中，一个点附近的函数写成 \(a/b\)，其中分母在该点不为零。这个条件只规定哪些元素可以作分母，并不依赖多项式的具体形式。现在从任意交换环出发，构造同时把一批元素变成单位的环。环中可能有零因子，因而需要保留第一章已经出现的现象：允许更多分母，有时也会使非零元素变为零。

\ProseParagraphBreak

本章的环均交换且含幺，环同态保持单位元，模满足 \(1m=m\)。允许零环；局部环则按非零环理解。除非另作说明，基域不要求代数闭。局部化的基本构造和理想对应可对照\cite[第4节，第20--24页]{Matsumura1989}。

\subsection{乘法集、分式及泛性质}
\label{g5:subsec:0201:01}

设 \(A\) 为环。包含 \(1\)、且对乘法封闭的子集 \(S\subseteq A\) 称为\term[chengfaji]{乘法集}[multiplicative set]。典型例子是
\[
	\{1,f,f^2,\ldots\},\qquad A\setminus\mathfrak p
	\quad(\mathfrak p\text{ 为素理想}).
\]
后一集合封闭正是素性条件的另一种说法。这里允许 \(0\in S\)；这种情形最终得到零环。

\ProseParagraphBreak

希望把 \((a,s)\in A\times S\) 写成分式 \(a/s\)。若允许分母成为单位，则任何被 \(S\) 中某个元素杀掉的元素也必须成为零。因此适当的等价关系是
\[
	(a,s)\sim(b,t)
	\quad\Longleftrightarrow\quad
	u(ta-sb)=0\text{ 对某个 }u\in S.
\]
自反性和对称性直接成立。若 \(u(ta-sb)=0\) 且 \(v(rb-tc)=0\)，则
\[
	uvt(ra-sc)
	=uv\bigl[r(ta-sb)+s(rb-tc)\bigr]=0,
\]
而 \(uvt\in S\)，所以关系也传递。

\begin{definition}\label{g5:def:ring-localization}
	设 \(A\) 为交换含幺环，\(S\subseteq A\) 为乘法集。以上等价类组成的集合记作 \(S^{-1}A\)，定义
	\[
	\frac as+\frac bt=\frac{at+bs}{st},\qquad
	\frac as\frac bt=\frac{ab}{st}.
	\]
	这个环称为 \(A\) 关于 \(S\) 的\term[jubuhua]{局部化}[localization]，自然同态为 \(\iota:A\to S^{-1}A,\ a\mapsto a/1\)。 对 \(S=\{f^n:n\geq0\}\) 记作 \(A_f\)。
\end{definition}
\symidx{\(S^{-1}A,\ A_f\)：环的局部化}

说明运算不依赖代表元。例如 \(u(s'a-sa')=0\) 时，把 \(a/s\) 换成 \(a'/s'\)，加法所得两个分子的交叉差是 \(t^2(s'a-sa')\)，乘法所得的交叉差是 \(bt(s'a-sa')\)，都被 \(u\) 杀掉。再对第二个分式作同样比较即可。环公理由通分后的环公理给出，零元、单位元分别是 \(0/1,1/1\)，而 \(s/1\) 的逆为 \(1/s\)。

\begin{proposition}\label{g5:prop:localization-zero}
	设 \(A\) 为交换含幺环，\(S\) 为乘法集。则
	\[
	\frac as=0\ \Longleftrightarrow\
	ta=0\text{ 对某个 }t\in S,\qquad
	\ker\iota=\{a\in A:ta=0\text{ 对某个 }t\in S\}.
	\]
	\(S^{-1}A\) 为零环当且仅当 \(0\in S\)。自然同态 \(\iota\) 单射当且仅当对每个 \(s\in S\)，乘法 \(a\mapsto sa\) 在 \(A\) 上单射。
\end{proposition}
\begin{proof}
	第一式是 \(a/s=0/1\) 的定义，核的公式随之得到。取 \(a=s=1\)，可知 \(1/1=0/1\) 当且仅当某个 \(t\in S\) 为零。最后，核为零等价于没有非零元素被 \(S\) 中任何元素杀掉，这正是所述乘法映射的单射性。
\end{proof}

例如在 \(A=k[x,y]/(xy)\) 中，\(\overline y\ne0\)，但在 \(A_{\overline x}\) 中 \(\overline y/1=0\)，因为 \(\overline x\,\overline y=0\)。事实上
\[
	A_{\overline x}\cong k[x,x^{-1}],
\]
同构把 \(x\) 送到 \(x\)，把 \(y\) 送到零。局部化在这里保留 \(x\ne0\) 的轴上部分，也同时去掉了另一条轴的函数信息。

\begin{theorem}[局部化的泛性质]\label{g5:thm:localization-universal}
	设 \(A,B\) 为交换含幺环，\(S\subseteq A\) 为乘法集，\(\varphi:A\to B\) 为环同态。若每个 \(\varphi(s)\)、\(s\in S\) 都是 \(B\) 的单位，则存在唯一环同态
	\[
	\widetilde\varphi:S^{-1}A\longrightarrow B,\qquad
	\widetilde\varphi(a/s)=\varphi(a)\varphi(s)^{-1},
	\]
	满足 \(\widetilde\varphi\circ\iota=\varphi\)。
\end{theorem}
\begin{proof}
	若 \(u(ta-sb)=0\)，应用 \(\varphi\) 后，\(\varphi(u),\varphi(s),\varphi(t)\) 都可逆，所以 \(\varphi(a)\varphi(s)^{-1}=\varphi(b)\varphi(t)^{-1}\)。故公式良定义；通分公式说明它保持加法、乘法及单位元。另一方面，每个分式都是 \((a/1)(s/1)^{-1}\)，所以任何满足要求的同态都必须采用这个公式，唯一性成立。
\end{proof}

泛性质使不同的构造可直接比较。例如先把 \(f\) 变成单位，再把 \(g\) 变成单位，与一次把 \(fg\) 变成单位得到同构的环：
\[
	(A_f)_{g/1}\cong A_{fg}.
\]
原因是 \(fg\) 可逆当且仅当 \(f,g\) 都可逆；两边的同态由泛性质唯一确定，复合在 \(A\) 上为恒等，因而在局部化上也为恒等。

\subsection{素理想处的局部化}
\label{g5:subsec:0201:02}

把指定点之外不消失的函数都作分母，会得到局部环。对一般环，“指定点”的代数条件由一个素理想承担。先说明分式属于扩张理想的判据。若 \(I\subseteq A\) 是理想，则由 \(\iota(I)\) 生成的理想等于
\[
	S^{-1}I=\{a/s:a\in I,\ s\in S\}.
\]
有限个这种分式可通分，分子仍在 \(I\) 中，故右侧确实是理想。对任意代表元有
\begin{equation}\label{g5:eq:localized-ideal-membership}
	\frac as\in S^{-1}I
	\quad\Longleftrightarrow\quad
	ta\in I\text{ 对某个 }t\in S.
\end{equation}
若 \(a/s=b/u\)、\(b\in I\)，取 \(v\in S\) 使 \(v(ua-sb)=0\)，则 \(vua=vsb\in I\)。反过来，\(ta\in I\) 时写成 \(a/s=ta/(ts)\) 即可。

\begin{theorem}\label{g5:thm:prime-localization}
	设 \(A\) 为交换含幺环，\(S\subseteq A\) 为乘法集。扩张与收缩给出保包含关系的双射
	\[
	\{\mathfrak p\subseteq A:\mathfrak p\text{ 素且 }\mathfrak p\cap S=\varnothing\}
	\ \longleftrightarrow\ \operatorname{Spec}(S^{-1}A),
	\qquad \mathfrak p\longmapsto S^{-1}\mathfrak p.
	\]
	这里 \(\operatorname{Spec}R\) 暂表示环 \(R\) 的全部素理想所成的集合。 对 \(A\) 的任一素理想 \(\mathfrak p\)，环 \(A_{\mathfrak p}=(A\setminus\mathfrak p)^{-1}A\) 是局部环，其唯一极大理想为 \(\mathfrak pA_{\mathfrak p}\)。
\end{theorem}
\begin{proof}
	若 \(\mathfrak p\cap S=\varnothing\)，由\eqref{g5:eq:localized-ideal-membership}及素性， \(a/s\in S^{-1}\mathfrak p\) 等价于 \(a\in\mathfrak p\)。 因此扩张是真理想，且两个分式的乘积属于它时，至少一个分式属于它。它为素理想，收缩恰为 \(\mathfrak p\)。

	若 \(\mathfrak q\subset S^{-1}A\) 素，收缩 \(\mathfrak p=\iota^{-1}(\mathfrak q)\) 素，且与 \(S\) 不交，因为单位不属于真理想。 又因 \(s/1\) 可逆，\(a/s\in\mathfrak q\) 当且仅当 \(a/1\in\mathfrak q\)，即 \(a\in\mathfrak p\)。故 \(\mathfrak q=S^{-1}\mathfrak p\)，两种操作互逆。

	取 \(S=A\setminus\mathfrak p\)。若 \(a/s\notin\mathfrak pA_{\mathfrak p}\)，则 \(a\notin\mathfrak p\)，其逆就是 \(s/a\)。故该理想之外全是单位，一切真理想都包含在它里面，它是唯一极大理想。
\end{proof}
\symidx{\(A_{\mathfrak p}\)：环在素理想 \(\mathfrak p\) 处的局部化}

这里沿用第一章的定义：有唯一极大理想的非零环称为局部环，常写为 \((R,\mathfrak m)\)。任一非零环的真理想都包含于极大理想：包含该理想的真理想组成偏序集，每条链的并仍为真理想，因为若并含 \(1\)，链中某个理想就含 \(1\)；极大性原理遂给出极大元。因此在局部环中，非单位恰组成 \(\mathfrak m\)。特别地，\(1-a\) 在 \(a\in\mathfrak m\) 时必为单位，否则 \(1=(1-a)+a\in\mathfrak m\)。

\begin{example}\label{g5:exm:localizing-axis}
	设 \(k\) 为域，\(A=k[x,y]/(xy)\)，以下用 \(x,y\) 表示剩余类。 在素理想 \(\mathfrak p=(x)\) 处，\(y\) 可逆，故 \(x=0\)；所有非零 \(k[y]\) 元素也成为单位，得到
	\[
	A_{\mathfrak p}\cong k(y).
	\]
	在 \(\mathfrak m=(x,y)\) 处，\(x,y\) 都非零。例如若 \(x/1=0\)，则有 \(s\notin\mathfrak m\) 使 \(sx\in(xy)\) 于 \(k[x,y]\)；约去 \(x\) 得 \(s\in(y)\)，矛盾。对 \(y\) 同理。因此 \(A_{\mathfrak m}\) 仍有两个非零元素之积为零，能够看到两条轴在原点相遇。
\end{example}

整环 \(A\) 在零素理想处的局部化是分式域 \(\operatorname{Frac}A\)。对素数 \(p\)，\(\mathbb Z_{(p)}\) 的元素是分母不被 \(p\) 整除的有理数，其极大理想由 \(p\) 生成。这些例子表明，局部化既可取在几何点处，也可取在一个不可约分支所对应的素理想处。

\subsection{单位、剩余域与局部同态}
\label{g5:subsec:0201:03}

局部环把“不在指定素理想中”变成“可除”。在一般乘法集下，单位也有一个直接的判据。

\begin{proposition}\label{g5:prop:units-residue-fields}
	设 \(A\) 为交换含幺环，\(S\) 为乘法集。分式 \(a/s\in S^{-1}A\) 可逆当且仅当 \((a)\cap S\ne\varnothing\)。 对每个素理想 \(\mathfrak p\subset A\)，有自然同构
	\[
	\kappa(\mathfrak p):=A_{\mathfrak p}/\mathfrak pA_{\mathfrak p}
	\cong\operatorname{Frac}(A/\mathfrak p).
	\]
\end{proposition}
\begin{proof}
	若 \(ab\in S\)，则 \(sb/(ab)\) 是 \(a/s\) 的逆。反过来，若 \((a/s)(b/t)=1\)，存在 \(u\in S\) 使 \(u(ab-st)=0\)。于是 \(a(ub)=ust\in S\)，得到判据。

	映射
	\[
	A_{\mathfrak p}\longrightarrow\operatorname{Frac}(A/\mathfrak p),
	\qquad a/s\longmapsto\overline a/\overline s
	\]
	良定义，因为 \(s\notin\mathfrak p\)。分式域中的每个分母都可如此提升，故映射满。其核恰为分子属于 \(\mathfrak p\) 的分式，即 \(\mathfrak pA_{\mathfrak p}\)，商环同构成立。
\end{proof}

对局部环 \((R,\mathfrak m)\)，域 \(R/\mathfrak m\) 称为它的\term[shengyuyu]{剩余域}[residue field]。上式的 \(\kappa(\mathfrak p)\) 是点 \(\mathfrak p\) 的剩余域。\symidx{\(\kappa(\mathfrak p)\)：素理想 \(\mathfrak p\) 处的剩余域} 在有限型代数闭域代数的极大理想处，弱零点定理使剩余域等于基域；在 \(k[x,y]/(xy)\) 的 \((x)\) 处，剩余域却是 \(k(y)\)。两者记录的分别是闭点上的值和沿整条不可约分支的有理函数值。

\begin{definition}\label{g5:def:local-homomorphism}
	设 \((R,\mathfrak m)\)、\((T,\mathfrak n)\) 为局部环。满足 \(\varphi(\mathfrak m)\subseteq\mathfrak n\) 的环同态 \(\varphi:R\to T\) 称为\term[jubutongtai]{局部同态}[local homomorphism]。
\end{definition}

这个条件等价于 \(\varphi^{-1}(\mathfrak n)=\mathfrak m\)：\(\mathfrak m\) 外的元素是单位，其像仍为单位，因而不可能落在 \(\mathfrak n\) 中。局部同态于是诱导剩余域的嵌入 \(R/\mathfrak m\hookrightarrow T/\mathfrak n\)。 一般的局部环之间的环同态未必局部。例如
\[
	k[t]_{(t)}\hookrightarrow k(t)
\]
把非单位 \(t\) 送成单位，靶的极大理想是零，所以不是局部同态。

\ProseParagraphBreak

若 \(\varphi:A\to B\) 为环同态，\(\mathfrak q\subset B\) 素，令 \(\mathfrak p=\varphi^{-1}(\mathfrak q)\)，则有自然局部同态
\[
	A_{\mathfrak p}\longrightarrow B_{\mathfrak q},
	\qquad a/s\longmapsto\varphi(a)/\varphi(s).
\]
分母可用，是因为 \(s\notin\mathfrak p\) 意味着 \(\varphi(s)\notin\mathfrak q\)；分子在 \(\mathfrak p\) 时像在 \(\mathfrak q\)，故同态局部。这是以后几何映射在点附近诱导环同态的基本形式。

\subsection{理想扩缩和饱和}
\label{g5:subsec:0201:04}

对环同态 \(\varphi:A\to B\)，理想 \(I\subseteq A\) 的\term[kuozhanglixiang]{扩张}[extension]是 \(IB\)，由 \(\varphi(I)\) 生成；理想 \(J\subseteq B\) 的\term[shousuolixiang]{收缩}[contraction]是 \(\varphi^{-1}(J)\)。一般只有 \(I\subseteq(IB)^c\)、\(J^cB\subseteq J\)。局部化时，后一包含总是等号，前一包含的差别恰由可用分母决定。

\begin{proposition}\label{g5:prop:localization-ideals}
	设 \(A\) 为交换含幺环，\(S\subseteq A\) 为乘法集，\(\iota:A\to S^{-1}A\) 为自然同态。则
	\[
	\iota^{-1}(S^{-1}I)=
	I:S^\infty:=\{a\in A:sa\in I\text{ 对某个 }s\in S\},
	\qquad
	J=S^{-1}\iota^{-1}(J).
	\]
	因而扩缩给出 \(S^{-1}A\) 的理想与 \(A\) 中满足 \(I:S^\infty=I\) 的理想之间的双射。 若 \(\overline S\) 是 \(S\) 在 \(A/I\) 中的像，则
	\[
	S^{-1}A/S^{-1}I\cong\overline S^{-1}(A/I).
	\]
	若 \(A\) 是 Noether 环，则 \(S^{-1}A\) 也是 Noether 环。
\end{proposition}
\begin{proof}
	第一个公式是\eqref{g5:eq:localized-ideal-membership}在分母为 \(1\) 时的情形。 若 \(a/s\in J\)，乘以单位 \(s/1\) 得 \(a/1\in J\)，所以 \(a\in\iota^{-1}(J)\)；反向包含来自理想性。这证明第二个公式，进而证明双射。 映射 \(a/s\mapsto\overline a/\overline s\) 满，其核是满足某个 \(ta\in I\) 的分式，故核为 \(S^{-1}I\)，商环公式得证。 最后若 \(A\) Noether，每个收缩理想由有限个 \(a_1,\ldots,a_r\) 生成，其扩张由 \(a_1/1,\ldots,a_r/1\) 生成；第二个公式说明这已涵盖局部化中的每个理想。
\end{proof}

当 \(S=\{f^n:n\geq0\}\) 时，上式就是第一章消去计算中的饱和 \(I:f^\infty\)。例如 \((xy):x^\infty=(y)\) 于 \(k[x,y]\)，因为可逆的 \(x\) 不再对方程 \(xy=0\) 贡献一个分支。

\ProseParagraphBreak

现在给全部素理想一个拓扑，使“在某点附近”成为统一的说法。环 \(A\) 的\term[supuxiang]{素谱}[prime spectrum]是集合 \(\operatorname{Spec}A\)，其闭集定义为
\[
	V(I)=\{\mathfrak p\in\operatorname{Spec}A:I\subseteq\mathfrak p\},
	\qquad I\subseteq A\text{ 为理想}.
\]
素性给出 \(V(IJ)=V(I)\cup V(J)\)，且 \(\bigcap_\lambda V(I_\lambda)=V(\sum_\lambda I_\lambda)\)； 配合 \(V(0)=\operatorname{Spec}A\)、\(V(A)=\varnothing\)，得到 Zariski 拓扑。 记 \(D(f)=\{\mathfrak p:f\notin\mathfrak p\}\)。 这里的点包括非极大素理想；\(V(I)\) 表示含 \(I\) 的素理想集，与第一章的经典零点集须按所在空间区分。
\symidx{\(\operatorname{Spec}A\)：环 \(A\) 的素谱}

\begin{proposition}\label{g5:prop:radical-prime-intersection}
	设 \(A\) 为交换含幺环，\(I\subseteq A\) 为理想。则
	\[
	\sqrt I=\bigcap_{\mathfrak p\supseteq I,\ \mathfrak p\text{ 素}}\mathfrak p,
	\]
	其中空族理想的交约定为 \(A\)。因而 \(V(I)=V(J)\) 当且仅当 \(\sqrt I=\sqrt J\)。
\end{proposition}
\begin{proof}
	若 \(a^n\in I\subseteq\mathfrak p\)，素性给出 \(a\in\mathfrak p\)，所以左侧包含于右侧。 若 \(a\notin\sqrt I\)，在 \(A/I\) 中 \(\overline a\) 的任何幂都非零；故 \((A/I)_{\overline a}\ne0\)。取其一个极大理想，依次收缩到 \(A/I\) 和 \(A\)，得到一个含 \(I\) 而不含 \(a\) 的素理想。这证明反向包含。 最后，同一个闭集给出同一族素理想的交，所以给出同一根式；反方向由素理想都是根式理想立即得到。
\end{proof}

\begin{proposition}\label{g5:prop:spec-basic-opens}
	设 \(A\) 为交换含幺环。主开集 \(D(f)\) 构成 \(\operatorname{Spec}A\) 的开基， \(D(f)\cap D(g)=D(fg)\)，且每个 \(D(f)\) 都是拟紧的，即任意开覆盖均有有限子覆盖。 局部化的素理想对应给出同胚
	\[
	\operatorname{Spec}A_f\cong D(f).
	\]
	特别地，\(\operatorname{Spec}A=D(1)\) 拟紧。
\end{proposition}
\begin{proof}
	对任何理想 \(I\)，有 \(\operatorname{Spec}A\setminus V(I)=\bigcup_{a\in I}D(a)\)。 素性又给出交集公式，所以主开集构成开基。

	局部化的素理想对应把 \(\operatorname{Spec}A_f\) 双射到不含 \(f\) 的素理想集。 在此双射下，主开集 \(D(a/f^n)\) 对应 \(D(a)\cap D(f)\)，因此双射及其逆都连续，得到同胚。

	最后先证任意环 \(R\) 的素谱拟紧。把一个开覆盖细化为主开覆盖 \(\{D(a_\lambda)\}\)。 若 \((a_\lambda)_\lambda\) 是真理想，它包含于某个极大理想，该极大理想就不属于任何 \(D(a_\lambda)\)，矛盾。 所以 \(1=\sum_{i=1}^n b_i a_{\lambda_i}\) 对有限个指标成立。 这些 \(D(a_{\lambda_i})\) 已覆盖素谱，因为素理想不可能同时含所有 \(a_{\lambda_i}\)。 取回它们所属的原覆盖成员便得有限子覆盖。 再对 \(R=A_f\) 应用这一结论与刚证的同胚，得到 \(D(f)\) 拟紧。
\end{proof}

一点 \(\mathfrak p\) 的闭包是 \(V(\mathfrak p)\)：每个包含该点的闭集都含这个集合。它在 \(\mathfrak p\) 极大时才是闭点；对整环，零素理想的闭包是整个素谱。这样的点称为所在不可约闭集的\term[fanxingdian]{泛型点}[generic point]。确实，若 \(V(\mathfrak p)=V(I)\cup V(J)\)，点 \(\mathfrak p\) 必在其中一边，而该边既为闭集就含其闭包 \(V(\mathfrak p)\)，所以此闭集不可约。用一个点代表整条分支，与在该点取分式域是相容的：两者都让我们考察分支的一般位置。

% 纲要条目排为编号小节标题，并进入目录与书签；说明段仍保留于纲要注释中。

% !TeX root = ../../Book5.tex
% 纲要标题：### 2.2 有限模与正合性
% 纲要位置：几何学计划纲要.md，第 4787 行。
% 正文状态：未撰写；本文件目前是框架。

% 纲要内容开始
% 
% * 有限生成、有限呈示及关系矩阵
% * 核、余核和正合列
% * 张量积、Hom 及标量变换
% * 局部化的正合性
% 
% 纲要内容结束

\section{有限模与正合性}
\label{g5:sec:0202}

待撰写

\subsection{有限生成、有限呈示及关系矩阵}
\label{g5:subsec:0202:01}
待撰写

\subsection{核、余核和正合列}
\label{g5:subsec:0202:02}
待撰写

\subsection{张量积、Hom 及标量变换}
\label{g5:subsec:0202:03}
待撰写

\subsection{局部化的正合性}
\label{g5:subsec:0202:04}
待撰写
% 纲要条目排为编号小节标题，并进入目录与书签；说明段仍保留于纲要注释中。

% !TeX root = ../../Book5.tex
% 纲要标题：### 2.3 Nakayama 与局部生成
% 纲要位置：几何学计划纲要.md，第 4794 行。
% 正文状态：未撰写；本文件目前是框架。

% 纲要内容开始
% 
% * 行列式技巧
% * 局部 Nakayama 引理
% * 极小生成元和剩余向量空间
% * 生成元提升及有限性假设
% 
% 纲要内容结束

\section{Nakayama 与局部生成}
\label{g5:sec:0203}

待撰写

\subsection{行列式技巧}
\label{g5:subsec:0203:01}
待撰写

\subsection{局部 Nakayama 引理}
\label{g5:subsec:0203:02}
待撰写

\subsection{极小生成元和剩余向量空间}
\label{g5:subsec:0203:03}
待撰写

\subsection{生成元提升及有限性假设}
\label{g5:subsec:0203:04}
待撰写
% 纲要条目排为编号小节标题，并进入目录与书签；说明段仍保留于纲要注释中。

% !TeX root = ../../Book5.tex
% 纲要标题：### 2.4 支集与关联素理想
% 纲要位置：几何学计划纲要.md，第 5095 行。
% 正文状态：初稿；尚未正式审核。

% 纲要内容开始
% 
% * 支集和零化子
% * 关联素理想的存在和有限性
% * 零因子及嵌入分支
% * 准素分解的所需版本
% 
% 纲要内容结束

\section{支集与关联素理想}
\label{g5:sec:0204}

局部化能够使一个非零模变为零，也能够使模中的某个非零元素消失。于是有两个相近而不同的问题：这个模在哪些素理想处仍然非零？它的哪些部分能够被一个素理想精确地零化？前一个问题产生支集，后一个问题产生关联素理想。关联素理想还会记录集中在较小闭集上的代数信息，使我们能进一步把一个理想分解成保留幂次的准素分量。

\ProseParagraphBreak

以下在 \(\operatorname{Spec}A\) 上讨论支集，点是素理想，闭集 \(V(I)\) 由包含 \(I\) 的素理想组成。支集的基本定义适用于任意交换含幺环；关联素理想的存在、有限性和准素分解中的有限性条件将在各个结果中写明。可对照\cite[第4节，第25--26页；第6节，第38--42页]{Matsumura1989}阅读。

\subsection{支集和零化子}
\label{g5:subsec:0204:01}
例如，对 \(A/(f)\) 在 \(\mathfrak p\) 处局部化，如果 \(f\notin\mathfrak p\)，那么 \(f\) 成为单位，而它在这个模上作用为零，所以整个局部化模为零。若想同时处理多个关系以及一般的有限模，就要把“哪些标量把它杀掉”整理成一个理想。

\begin{definition}\label{g5:def:support-annihilator}
	设 \(A\) 为交换含幺环，\(M\) 为 \(A\) 模。元素 \(m\in M\) 的\term[linghuazi]{零化子}[annihilator]及模 \(M\) 的零化子分别是
	\[
		\operatorname{Ann}_A(m)=\{a\in A:am=0\},\qquad
		\operatorname{Ann}_A(M)=\{a\in A:aM=0\}
		=\bigcap_{m\in M}\operatorname{Ann}_A(m).
	\]
	它们都是 \(A\) 的理想。模 \(M\) 的\term[zhiji]{支集}[support]定义为
	\[
		\operatorname{Supp}_A(M)=\{\mathfrak p\in\operatorname{Spec}A:M_{\mathfrak p}\ne0\}.
	\]
	底环明确时省略下标 \(A\)。规定 \(\operatorname{Ann}_A(0)=A\)、\(\operatorname{Supp}_A(0)=\varnothing\)。
\end{definition}
\symidx{\(\operatorname{Ann}_A(M)\)：模的零化子}
\symidx{\(\operatorname{Supp}_A(M)\)：模的支集}

映射 \(A\to Am\)、\(a\mapsto am\) 的核是 \(\operatorname{Ann}_A(m)\)，所以
\[
	Am\cong A/\operatorname{Ann}_A(m).
\]
一个元素生成的循环子模，由它的零化子完全决定。对任意模则有 \(\operatorname{Ann}_A(M)\subseteq\operatorname{Ann}_A(m)\)；整个模的共同零化条件可能比任何一个元素的零化条件都少。

\begin{proposition}\label{g5:prop:support-annihilator}
	设 \(A\) 为交换含幺环，\(M\) 为 \(A\) 模。总有
	\[
		\operatorname{Supp}_A(M)\subseteq V(\operatorname{Ann}_A(M)).
	\]
	若 \(M\) 有限生成，则等号成立，支集因而是闭集。特别地，对任意理想 \(I\subseteq A\)，有 \(\operatorname{Supp}_A(A/I)=V(I)\)。
\end{proposition}
\begin{proof}
	若某个 \(a\in\operatorname{Ann}_A(M)\) 不属于 \(\mathfrak p\)，则 \(a/1\) 在 \(A_{\mathfrak p}\) 中可逆，却把 \(M_{\mathfrak p}\) 的每个元素送到零，所以 \(M_{\mathfrak p}=0\)。这证明一般的包含关系。

	现在设 \(M\) 由 \(m_1,\ldots,m_r\) 生成，且 \(M_{\mathfrak p}=0\)。由模局部化的零元判据\eqref{g5:eq:module-localization-zero}，对每个 \(i\) 存在 \(s_i\notin\mathfrak p\) 使 \(s_im_i=0\)。令 \(s=\prod_i s_i\)，则 \(s\notin\mathfrak p\)，且 \(s\) 杀掉全部生成元，从而 \(sM=0\)。故 \(\operatorname{Ann}_A(M)\) 不包含在 \(\mathfrak p\) 中。这给出有限生成情形的反向包含。最后，\(A/I\) 由 \(1+I\) 生成，且其零化子为 \(I\)，于是得到商模公式。
\end{proof}

证明中有限生成性的作用，是把各个元素所需的分母合成一个共同分母。对无限多个生成元，不能作这样的有限乘积。

\begin{example}\label{g5:exm:nonclosed-module-support}
	设 \(k\) 为代数闭域，\(A=k[t]\)，取无限直和
	\[
		M=\bigoplus_{a\in k}A/(t-a).
	\]
	一个多项式若杀掉整个 \(M\)，就在所有 \(a\in k\) 处取零，故 \(\operatorname{Ann}_A(M)=0\)。在 \((0)\) 处局部化时，每个直和分量都被可逆的 \(t-a\) 杀掉；直和中的每个元素只有有限个非零坐标，所以 \(M_{(0)}=0\)。在 \((t-a)\) 处，第 \(a\) 个分量仍非零。由于 \(k[t]\) 的非零素理想正是这些极大理想，
	\[
		\operatorname{Supp}_A(M)=\{(t-a):a\in k\},
		\qquad V(\operatorname{Ann}_A(M))=\operatorname{Spec}k[t].
	\]
	这个支集包含全部闭点而遗漏泛型点 \((0)\)。任何包含它的闭集 \(V(J)\) 都要求 \(J\) 中每个多项式在全部 \(a\) 处取零，因而 \(J=0\)。所以支集稠密，却不是闭集。
\end{example}

\begin{proposition}\label{g5:prop:support-exact}
	设 \(A\) 为交换含幺环，\(0\to L\to M\to N\to0\) 为 \(A\) 模的短正合列。则
	\[
		\operatorname{Supp}_A(M)=\operatorname{Supp}_A(L)\cup\operatorname{Supp}_A(N).
	\]
	这里不要求模有限生成。
\end{proposition}
\begin{proof}
	对任意 \(\mathfrak p\in\operatorname{Spec}A\)，\cref{g5:thm:localization-exact}给出短正合列
	\[
		0\longrightarrow L_{\mathfrak p}\longrightarrow M_{\mathfrak p}
		\longrightarrow N_{\mathfrak p}\longrightarrow0.
	\]
	中间项为零当且仅当两端都为零。对全部 \(\mathfrak p\) 应用这个等价关系，就得到支集公式。
\end{proof}

因此，通过一条有限滤过把模拆成较简单的商模时，原模的支集就是这些商模支集的并。接下来要构造一种商模形如 \(A/\mathfrak p\) 的滤过，使其支集成为素理想确定的不可约闭集。

\subsection{关联素理想的存在和有限性}
\label{g5:subsec:0204:02}
若一个非零元素 \(m\) 的零化子恰为素理想 \(\mathfrak p\)，则循环子模 \(Am\cong A/\mathfrak p\) 在整条 \(V(\mathfrak p)\) 上存活。它不仅指出了某个局部化非零，还在原模中找到了一个承载这条闭集的具体子模。

\begin{definition}\label{g5:def:associated-prime}
	设 \(A\) 为交换含幺环，\(M\) 为 \(A\) 模。若素理想 \(\mathfrak p\) 可以写成 \(\operatorname{Ann}_A(m)\)，其中 \(0\ne m\in M\)，则称 \(\mathfrak p\) 为 \(M\) 的\term[guanliansulixiang]{关联素理想}[associated prime ideal]。全部关联素理想组成的集合记为 \(\operatorname{Ass}_A(M)\)。等价地，\(\mathfrak p\in\operatorname{Ass}_A(M)\) 当且仅当存在单射 \(A/\mathfrak p\hookrightarrow M\)。零模的关联素理想集为空。
\end{definition}
\symidx{\(\operatorname{Ass}_A(M)\)：模的关联素理想集}

一般元素的零化子未必素。例如在 \(k[t]/(t^2)\) 中，\(1\) 的零化子为 \((t^2)\)，而 \(t\) 的零化子是 \((t)\)。把元素换成一个适当的非零倍式，可能增大它的零化子；Noether 条件保证这个过程不会无限严格上升。

\begin{lemma}\label{g5:lem:associated-prime-exists}
	设 \(A\) 为交换含幺 Noether 环，\(M\ne0\) 为任意 \(A\) 模。在集合
	\[
		\{\operatorname{Ann}_A(m):0\ne m\in M\}
	\]
	中按包含关系取的每个极大元都是素理想。因此 \(\operatorname{Ass}_A(M)\ne\varnothing\)。
\end{lemma}
\begin{proof}
	Noether 环的理想升链条件保证这个非空理想族有极大元。取 \(\mathfrak p=\operatorname{Ann}_A(m)\) 为其中一个；由于 \(m\ne0\)，有 \(\mathfrak p\ne A\)。若 \(ab\in\mathfrak p\)、\(b\notin\mathfrak p\)，则 \(bm\ne0\)，并有
	\[
		\mathfrak p\subseteq\operatorname{Ann}_A(bm).
	\]
	极大性迫使二者相等。由 \(abm=0\)，得到 \(a\in\operatorname{Ann}_A(bm)=\mathfrak p\)，所以 \(\mathfrak p\) 为素理想。
\end{proof}

这个引理只要求环 Noether，没有要求 \(M\) 有限生成。有限生成性将用于控制关联素理想的总数。为此先说明，关联素理想怎样经过短正合列。

\begin{proposition}\label{g5:prop:ass-exact-sequence}
	设 \(A\) 为交换含幺环，\(0\to L\to M\to N\to0\) 为短正合列，则
	\[
		\operatorname{Ass}_A(L)\subseteq\operatorname{Ass}_A(M)
		\subseteq\operatorname{Ass}_A(L)\cup\operatorname{Ass}_A(N).
	\]
	特别地，有限个模的直和满足
	\[
		\operatorname{Ass}_A(M_1\oplus\cdots\oplus M_r)
		=\bigcup_{i=1}^r\operatorname{Ass}_A(M_i).
	\]
\end{proposition}
\begin{proof}
	将 \(L\) 识别为 \(M\) 的子模。其非零元素的零化子不变，所以第一个包含成立。对 \(\mathfrak p\in\operatorname{Ass}_A(M)\)，取循环子模 \(C=Am\cong A/\mathfrak p\)。若 \(C\cap L=0\)，则 \(C\to N\) 为单射，故 \(\mathfrak p\in\operatorname{Ass}_A(N)\)。

	若 \(C\cap L\ne0\)，选 \(0\ne bm\in C\cap L\)。由于 \(\operatorname{Ann}_A(m)=\mathfrak p\)，有 \(b\notin\mathfrak p\)。对于 \(a\in A\)，
	\[
		abm=0\quad\Longleftrightarrow\quad ab\in\mathfrak p
		\quad\Longleftrightarrow\quad a\in\mathfrak p,
	\]
	最后一个等价使用了素性及 \(b\notin\mathfrak p\)。因此 \(\operatorname{Ann}_A(bm)=\mathfrak p\)，说明 \(\mathfrak p\in\operatorname{Ass}_A(L)\)。这证明第二个包含。对直和，短正合列给出向右的包含，每个直和分量嵌入整个直和又给出反向包含，有限次归纳即得公式。
\end{proof}

关联素理想在商模中的行为与子模不同。例如取域 \(k\) 上的 \(A=k[t]\)，短正合列
\[
	0\longrightarrow A\xrightarrow{\,t\,}A\longrightarrow A/(t)\longrightarrow0
\]
中，前两个非零模的关联素理想均为 \((0)\)，商模却出现 \((t)\)。因此不能把上个命题中的包含关系一概改成等号。

\begin{theorem}\label{g5:thm:prime-filtration}
	设 \(A\) 为交换含幺 Noether 环，\(M\) 为有限生成 \(A\) 模。则存在有限滤过
	\[
		0=M_0\subsetneq M_1\subsetneq\cdots\subsetneq M_r=M,
		\qquad M_i/M_{i-1}\cong A/\mathfrak p_i,
	\]
	其中各 \(\mathfrak p_i\) 为素理想；\(M=0\) 时允许 \(r=0\)。这样的滤过称为\term[sulvguo]{素滤过}[prime filtration]，并满足
	\[
		\operatorname{Ass}_A(M)\subseteq\{\mathfrak p_1,\ldots,\mathfrak p_r\},
		\qquad \operatorname{Supp}_A(M)=\bigcup_{i=1}^r V(\mathfrak p_i).
	\]
	特别地，\(\operatorname{Ass}_A(M)\) 为有限集合。
\end{theorem}
\begin{proof}
	若 \(M\ne0\)，由\cref{g5:lem:associated-prime-exists}选择 \(M_1\cong A/\mathfrak p_1\) 的非零循环子模。若 \(M_1\ne M\)，对非零商模 \(M/M_1\) 作同样选择，再把选出的循环子模拉回 \(M\)，得到 \(M_2\)。如此继续，每一步都严格增大子模。由\cref{g5:prop:noetherian-finite-modules}，有限生成 \(A\) 模满足子模升链条件，因此只能作有限步，并且终止时必须达到 \(M\)。

	对素理想 \(\mathfrak p\)，模 \(A/\mathfrak p\) 中每个非零元素的零化子都是 \(\mathfrak p\)，所以 \(\operatorname{Ass}_A(A/\mathfrak p)=\{\mathfrak p\}\)。对滤过中逐个短正合列使用\cref{g5:prop:ass-exact-sequence}，得到关联素理想的包含关系。同样使用\cref{g5:prop:support-exact,g5:prop:support-annihilator}，得到支集的有限并公式。
\end{proof}

素滤过中的 \(\mathfrak p_i\) 可能重复，也可能包含并非 \(M\) 的关联素理想的素理想；滤过是构造工具，\(\operatorname{Ass}_A(M)\) 才是由模本身确定的集合。现在考察局部化怎样从这个集合中选出仍然可见的素理想。

\begin{theorem}\label{g5:thm:associated-localization}
	设 \(A\) 为交换含幺 Noether 环，\(M\) 为 \(A\) 模，\(S\subseteq A\) 为乘法集。则
	\[
		\operatorname{Ass}_{S^{-1}A}(S^{-1}M)
		=\{S^{-1}\mathfrak p:\mathfrak p\in\operatorname{Ass}_A(M),\quad
		\mathfrak p\cap S=\varnothing\}.
	\]
	左右的素理想通过\cref{g5:thm:prime-localization}的扩张、收缩对应。
\end{theorem}
\begin{proof}
	若 \(\mathfrak p=\operatorname{Ann}_A(m)\) 且 \(\mathfrak p\cap S=\varnothing\)，则 \(m/1\ne0\)，否则某个 \(s\in S\) 满足 \(sm=0\)，与不交条件矛盾。又有
	\[
		\frac as\frac m1=0
		\quad\Longleftrightarrow\quad tam=0\text{ 对某个 }t\in S
		\quad\Longleftrightarrow\quad a\in\mathfrak p.
	\]
	最后一步使用 \(ta\in\mathfrak p\)、\(t\notin\mathfrak p\) 及素性。因此 \(m/1\) 的零化子恰为 \(S^{-1}\mathfrak p\)。

	反过来，设 \(\mathfrak q\) 是某个非零 \(m/s\in S^{-1}M\) 的零化子。因为 \(s/1\) 可逆，\(m/1\) 的零化子同样是 \(\mathfrak q\)。令 \(\mathfrak p\subseteq A\) 为 \(\mathfrak q\) 的收缩；它为素理想、与 \(S\) 不交，且 \(\mathfrak q=S^{-1}\mathfrak p\)。Noether 性使 \(\mathfrak p=(a_1,\ldots,a_\ell)\) 有限生成。每个 \(a_i\) 在局部化后杀掉 \(m/1\)，所以存在 \(t_i\in S\) 使 \(t_i a_im=0\)。取 \(t=\prod_i t_i\)，便有 \(\mathfrak p\cdot tm=0\)。

	元素 \(tm\) 非零，因为其局部化是单位 \(t/1\) 乘以非零的 \(m/1\)。若 \(b\cdot tm=0\)，局部化并约去这个单位，得到 \((b/1)(m/1)=0\)，故 \(b\in\mathfrak p\)。于是 \(\operatorname{Ann}_A(tm)=\mathfrak p\)，证明 \(\mathfrak p\in\operatorname{Ass}_A(M)\)。若 \(0\in S\)，局部化环和模都为零，右侧也没有与 \(S\) 不交的素理想，公式仍成立。
\end{proof}

这个反向证明解释了为何需要把局部化中的元素乘以一个共同分母：最初的 \(m\) 未必已被整个 \(\mathfrak p\) 杀掉，而 \(tm\) 可以做到。所用有限性来自素理想 \(\mathfrak p\) 的有限生成，因此这里仍不需要 \(M\) 有限生成。

\subsection{零因子及嵌入分支}
\label{g5:subsec:0204:03}
支集知道模在哪里存在，却没有直接告诉我们哪些标量的乘法会失去单射性。关联素理想恰好补足这一信息：它们的并就是所有能杀掉非零元素的标量。

\begin{definition}\label{g5:def:module-zero-divisor}
	设 \(A\) 为交换含幺环，\(M\) 为 \(A\) 模。若存在 \(0\ne m\in M\) 使 \(am=0\)，则称 \(a\in A\) 为 \(M\) 上的\term[lingyinzi]{零因子}[zero divisor]。其余元素的乘法映射 \(M\xrightarrow{a}M\) 为单射，称为 \(M\) 上的\term[zhengzeyuan]{正则元}[regular element]。特别地，\(M\ne0\) 时零元素也是零因子；零模上每个元素都是正则元。
\end{definition}

\begin{theorem}\label{g5:thm:zero-divisors-associated}
	设 \(A\) 为交换含幺 Noether 环，\(M\) 为任意 \(A\) 模。则 \(M\) 上的零因子集恰为
	\[
		\bigcup_{\mathfrak p\in\operatorname{Ass}_A(M)}\mathfrak p.
	\]
\end{theorem}
\begin{proof}
	若 \(a\in\mathfrak p=\operatorname{Ann}_A(m)\)，其中 \(m\ne0\)，则 \(am=0\)，所以 \(a\) 为零因子。反过来，若 \(a\) 为零因子，则子模 \(L=\{m\in M:am=0\}\) 非零。由\cref{g5:lem:associated-prime-exists}，\(L\) 中有一个非零元素 \(m\)，其零化子是素理想 \(\mathfrak p\)。这个零化子在 \(M\) 中仍相同，所以 \(\mathfrak p\in\operatorname{Ass}_A(M)\)，而 \(am=0\) 又给出 \(a\in\mathfrak p\)。
\end{proof}

在 Noether 环的有限模上，检验一个元素是否正则因而只需避开有限多个关联素理想。这种判定保留了支集中可能看不见的较小部分；先比较两个集合各自的极小元。

\begin{corollary}\label{g5:cor:minimal-support-associated}
	设 \(A\) 为交换含幺 Noether 环，\(M\) 为有限生成 \(A\) 模。则
	\[
		\operatorname{Ass}_A(M)\subseteq\operatorname{Supp}_A(M),\qquad
		\operatorname{Supp}_A(M)=
		\bigcup_{\mathfrak p\in\operatorname{Ass}_A(M)}V(\mathfrak p).
	\]
	按理想包含关系，\(\operatorname{Supp}_A(M)\) 与 \(\operatorname{Ass}_A(M)\) 有相同的极小元。它们为有限多个，其闭包恰是支集的不可约分支。
\end{corollary}
\begin{proof}
	若 \(\mathfrak p=\operatorname{Ann}_A(m)\)、\(m\ne0\)，则对每个包含 \(\mathfrak p\) 的素理想 \(\mathfrak q\)，元素 \(m/1\in M_{\mathfrak q}\) 非零：否则有 \(s\notin\mathfrak q\) 使 \(sm=0\)，就会有 \(s\in\mathfrak p\subseteq\mathfrak q\)，矛盾。因此 \(V(\mathfrak p)\subseteq\operatorname{Supp}_A(M)\)，也得到第一个包含。

	现在取 \(\mathfrak q\in\operatorname{Supp}_A(M)\)。由\cref{g5:prop:localization-ideals}，\(A_{\mathfrak q}\) 为 Noether 环，而 \(M_{\mathfrak q}\ne0\)。\cref{g5:lem:associated-prime-exists}保证它有一个关联素理想。再由\cref{g5:thm:associated-localization}，该素理想来自某个 \(\mathfrak p\in\operatorname{Ass}_A(M)\)，且 \(\mathfrak p\subseteq\mathfrak q\)。于是 \(\mathfrak q\in V(\mathfrak p)\)，证明并集公式。

	由于关联素理想集有限，其中每个元素都包含一个极小元。因此上式可以只对这些极小元取并，而任意支集素理想都包含其中一个。结合 \(\operatorname{Ass}_A(M)\subseteq\operatorname{Supp}_A(M)\)，两个集合的极小元恰好相同。每个 \(V(\mathfrak p)\) 都不可约，且理想之间互不包含使这些闭集也互不包含。有限个闭集的并中的任一不可约子集必落在其中一项内，否则把相交拆开就违反不可约性。故这些 \(V(\mathfrak p)\) 正是支集的不可约分支。
\end{proof}

对 Noether 环上的有限模，支集中的极小素理想称为\term[guliguanliansulixiang]{孤立关联素理想}[isolated associated prime]，也称极小关联素理想；其余关联素理想称为\term[qianrusulixiang]{嵌入素理想}[embedded prime]。若 \(\mathfrak p\) 为嵌入素理想，它包含一个更小的关联素理想 \(\mathfrak q\)，所以 \(V(\mathfrak p)\) 严格包含在 \(V(\mathfrak q)\) 中。它所记录的子模集中在已有不可约分支的一个较小闭集上。

\begin{example}\label{g5:exm:associated-embedded-point}
	设 \(k\) 为任意域，\(R=k[x,y]\)，\(I=(x^2,xy)\)，\(M=R/I\)。这是\cref{g5:ex:radical-embedded-model}中已出现的模型；现在把它作为 \(R\) 模计算。以 \(\overline x,\overline y\) 表示剩余类，每个元素唯一写成
	\[
		a(\overline y)+b\overline x,\qquad a(y)\in k[y],\quad b\in k.
	\]
	存在性来自 \(\overline x^2=\overline x\,\overline y=0\)；唯一性来自 \(I\) 中每个单项式都含 \(x^2\) 或 \(xy\)，不能出现非零的 \(a(y)+bx\)。

	因为 \(f(x,y)\overline x=f(0,0)\overline x\)，且 \(\overline x\ne0\)，有
	\[
		\operatorname{Ann}_R(\overline x)=(x,y).
	\]
	另一方面，\(f(x,y)\overline y=f(0,y)\overline y\)，而 \(k[y]\) 是整环，所以
	\[
		\operatorname{Ann}_R(\overline y)=(x).
	\]
	两者都是关联素理想。为了排除其他关联素理想，取子模 \(R\overline x\cong R/(x,y)\)，其商模为 \(R/(x)\)。短正合列
	\[
		0\longrightarrow R/(x,y)\xrightarrow{\,1\mapsto\overline x\,}M
		\longrightarrow R/(x)\longrightarrow0
	\]
	与\cref{g5:prop:ass-exact-sequence}给出
	\[
		\operatorname{Ass}_R(M)=\{(x),(x,y)\},\qquad
		\operatorname{Supp}_R(M)=V(x).
	\]
	支集只有一条不可约分支 \(V(x)\)，但 \(R\overline x\) 是支撑在原点 \((x,y)\) 上的非零子模。其元素 \(\overline x\) 在 \(y\) 可逆时消失，在原点局部化后仍非零，因为没有原点之外的分母能落入它的零化子 \((x,y)\)。因此 \((x)\) 为孤立关联素理想，\((x,y)\) 为嵌入素理想。零因子集为 \((x)\cup(x,y)=(x,y)\)；例如 \(y\) 是零因子，虽然它不属于唯一的极小关联素理想 \((x)\)。
\end{example}

只检查支集的不可约分支，会漏掉上例中 \(y\) 杀掉 \(\overline x\) 的现象。准素分解的目标就是在分解方程理想时保留这种信息，同时区分哪些分量由理想唯一决定。

\subsection{准素分解的所需版本}
\label{g5:subsec:0204:04}
第一章的根式理想与不可约分支对应，使我们能把代数集拆成不可约闭集。但取根式会丢掉 \(x^2\) 与 \(x\) 的区别；本小节直接分解理想，允许一个分量内部带有幂零元素。与素理想相比，适当放宽的条件是：乘积为零时，一个因子可能只在取幂后才变为零。

\begin{definition}\label{g5:def:primary-ideal}
	设 \(A\) 为交换含幺环。真理想 \(Q\subsetneq A\) 若满足
	\[
		ab\in Q,\quad b\notin Q
		\quad\Longrightarrow\quad a^n\in Q\text{ 对某个 }n\ge1,
	\]
	则称为\term[zhunsulixiang]{准素理想}[primary ideal]。若 \(\sqrt Q=\mathfrak p\)，则称 \(Q\) 为 \(\mathfrak p\) 准素理想。把真理想写成有限个准素理想的交，称为\term[zhunsufenjie]{准素分解}[primary decomposition]。
\end{definition}

等价地，\(Q\) 为准素理想，当且仅当商环 \(A/Q\) 的每个零因子都是幂零元。定义中的 \(b\notin Q\) 正保证 \(b+Q\) 是检验零因子的非零元素。素理想当然准素，此时可以取 \(n=1\)。

\begin{lemma}\label{g5:lem:primary-ideal-properties}
	设 \(A\) 为交换含幺环，\(Q\subsetneq A\) 为准素理想，\(\mathfrak p=\sqrt Q\)。则 \(\mathfrak p\) 是素理想。对 \(c\notin Q\)，理想
	\[
		Q:c=\{a\in A:ac\in Q\}
	\]
	仍为 \(\mathfrak p\) 准素理想。有限个 \(\mathfrak p\) 准素理想的交也为 \(\mathfrak p\) 准素理想。此外，若 \(\iota:A\to A_{\mathfrak p}\) 为自然同态，则
	\[
		\iota^{-1}(QA_{\mathfrak p})=Q.
	\]
\end{lemma}
\begin{proof}
	若 \(ab\in\sqrt Q\)、\(b\notin\sqrt Q\)，则对某个 \(r\) 有 \(a^rb^r\in Q\)，而 \(b^r\notin Q\)。准素性给出 \(a^{rn}\in Q\)，所以 \(a\in\sqrt Q\)。这证明 \(\mathfrak p\) 为素理想。

	由于 \(c\notin Q\)，\(Q:c\) 是真理想。若 \(ab\in Q:c\)、\(b\notin Q:c\)，则 \(abc\in Q\)、\(bc\notin Q\)，故某个 \(a^n\in Q\subseteq Q:c\)。所以 \(Q:c\) 准素。包含关系 \(Q\subseteq Q:c\) 给出 \(\mathfrak p\subseteq\sqrt{Q:c}\)；反之，若 \(a^r c\in Q\)，由 \(c\notin Q\) 得到 \(a\in\mathfrak p\)，所以两根式相同。

	设 \(Q_1,\ldots,Q_r\) 都是 \(\mathfrak p\) 准素理想。有限交的根式仍为 \(\mathfrak p\)：一个元素分别有幂属于各个 \(Q_i\)，取指数的最大值即可使同一个幂属于交。若 \(ab\in\bigcap_iQ_i\)、\(b\notin\bigcap_iQ_i\)，选择一个 \(b\notin Q_j\)，则 \(a\in\mathfrak p\)，进而某个 \(a^N\in\bigcap_iQ_i\)。因此该交准素。

	最后，若 \(a/1\in QA_{\mathfrak p}\)，由\cref{g5:prop:localization-ideals}，存在 \(s\notin\mathfrak p\) 使 \(sa\in Q\)。如果 \(a\notin Q\)，准素性将给出 \(s\in\sqrt Q=\mathfrak p\)，矛盾。故 \(a\in Q\)，证明收缩公式。
\end{proof}

准素性也可以完全用关联素理想表述。

\begin{proposition}\label{g5:prop:primary-associated-singleton}
	设 \(A\) 为交换含幺 Noether 环，\(Q\subsetneq A\) 为理想。则 \(Q\) 准素，当且仅当 \(\operatorname{Ass}_A(A/Q)\) 只有一个元素。在这种情形下，
	\[
		\operatorname{Ass}_A(A/Q)=\{\sqrt Q\}.
	\]
\end{proposition}
\begin{proof}
	若 \(Q\) 准素，令 \(\mathfrak p=\sqrt Q\)。取任意 \(\mathfrak q\in\operatorname{Ass}_A(A/Q)\)，它是某个非零剩余类的零化子。由于 \(Q\) 杀掉全部商模，有 \(Q\subseteq\mathfrak q\)，故 \(\mathfrak p\subseteq\mathfrak q\)。另一方面，\(\mathfrak q\) 中每个元素都在 \(A/Q\) 上是零因子，准素性使它模 \(Q\) 幂零，所以 \(\mathfrak q\subseteq\mathfrak p\)。关联素理想的存在性保证集合非空，于是它恰为 \(\{\mathfrak p\}\)。

	反过来，若 \(\operatorname{Ass}_A(A/Q)=\{\mathfrak p\}\)，则\cref{g5:cor:minimal-support-associated,g5:prop:support-annihilator}给出 \(V(Q)=V(\mathfrak p)\)。由\cref{g5:prop:radical-prime-intersection}，\(\sqrt Q=\mathfrak p\)。\cref{g5:thm:zero-divisors-associated}又说明 \(A/Q\) 上的零因子恰为 \(\mathfrak p\) 中的元素，它们模 \(Q\) 都幂零，故 \(Q\) 准素。
\end{proof}

为了证明每个理想都有这样的分解，先作一个较粗的拆分。一个真理想 \(J\) 若不能表示成两个严格包含它的理想之交，就称为\term[bukeyuelixiang]{不可约理想}[irreducible ideal]。这里不可约描述的是理想的交分解，与闭集的不可约性需要分别判断；后者只看根式。

\begin{lemma}\label{g5:lem:irreducible-ideal-primary}
	设 \(A\) 为交换含幺 Noether 环。每个真理想都是有限个不可约理想的交，而且每个不可约理想均准素。
\end{lemma}
\begin{proof}
	先证明有限分解。若存在不能作有限不可约分解的真理想，则由理想升链条件，可在这些反例中取一个按包含关系极大的 \(J\)。它本身不可能不可约，故有 \(J=J_1\cap J_2\)，其中两个 \(J_i\) 都严格包含 \(J\)。两者仍为真理想，否则这个交就等于另一项而无法严格增大。由 \(J\) 的极大性，每个 \(J_i\) 都是有限个不可约理想的交，将两份表达式合并，就给出 \(J\) 的有限不可约分解，矛盾。

	再设 \(Q\) 不可约。转到商环 \(B=A/Q\)，其中零理想不可约。若 \(ab=0\) 且 \(b\ne0\)，要证明 \(a\) 幂零。考虑升链
	\[
		(0:a)\subseteq(0:a^2)\subseteq\cdots,
		\qquad (0:a^n)=\{u\in B:a^nu=0\}.
	\]
	商环 \(B\) 的理想升链提升为 \(A\) 中含 \(Q\) 的理想升链，故 \(B\) 仍 Noether。所以可以取 \(N\ge1\) 使 \((0:a^N)=(0:a^{N+1})\)。现在若 \(z\in(a^N)\cap(b)\)，写成 \(z=a^Nu=bv\)。由 \(ab=0\)，有 \(a^{N+1}u=az=0\)；稳定性遂给出 \(a^Nu=0\)，也就是 \(z=0\)。因此
	\[
		(0)=(a^N)\cap(b).
	\]
	零理想不可约而 \((b)\ne0\)，迫使 \((a^N)=0\)。这就证明 \(a^N=0\)，从而 \(Q\) 准素。
\end{proof}

\begin{theorem}[准素分解定理]\label{g5:thm:primary-decomposition}
	\termidx[zhunsufenjiedingli]{准素分解定理}\foreignidx{primary decomposition theorem}设 \(A\) 为交换含幺 Noether 环，\(I\subsetneq A\) 为真理想。则存在有限分解
	\[
		I=Q_1\cap\cdots\cap Q_r,
	\]
	其中每个 \(Q_i\) 均准素。可以要求这份分解不可删去任何一项，且 \(\mathfrak p_i=\sqrt{Q_i}\) 两两不同。满足这两个条件的分解称为\term[jixiaozhunsufenjie]{极小准素分解}[minimal primary decomposition]。
\end{theorem}
\begin{proof}
	由\cref{g5:lem:irreducible-ideal-primary}，先将 \(I\) 写成有限个不可约理想的交，它们都是准素理想。将根式相同的项合并，由\cref{g5:lem:primary-ideal-properties}，合并后的交仍是同一个素理想对应的准素理想。再删去不影响交的冗余项，有限次后便得到所需分解。
\end{proof}

这个定理保证分解存在。分解中的根式与实际准素理想承担不同信息：根式指出分量所在的闭集，准素理想还保留闭集上的幂次。下一定理分别判断它们的唯一性；对应结论可参见\cite[第6节，定理6.8]{Matsumura1989}。

\begin{theorem}\label{g5:thm:primary-uniqueness}
	设 \(A\) 为交换含幺 Noether 环，\(I\subsetneq A\)，并取一个极小准素分解
	\[
		I=\bigcap_{i=1}^rQ_i,\qquad \mathfrak p_i=\sqrt{Q_i}.
	\]
	则
	\[
		\operatorname{Ass}_A(A/I)=\{\mathfrak p_1,\ldots,\mathfrak p_r\}.
	\]
	因此分解中的素理想集与分量数由 \(I\) 唯一决定。若 \(\mathfrak p_i\) 为其中的极小元，则对应的孤立准素分量也唯一，且
	\[
		Q_i=\{a\in A:a/1\in IA_{\mathfrak p_i}\}.
	\]
\end{theorem}
\begin{proof}
	自然的对角映射
	\[
		A/I\longrightarrow\bigoplus_{i=1}^r A/Q_i,\qquad
		a+I\longmapsto(a+Q_1,\ldots,a+Q_r)
	\]
	为单射，因为它的核由 \(\bigcap_iQ_i=I\) 决定。由\cref{g5:prop:ass-exact-sequence,g5:prop:primary-associated-singleton}，得到
	\[
		\operatorname{Ass}_A(A/I)\subseteq\bigcup_i\operatorname{Ass}_A(A/Q_i)
		=\{\mathfrak p_1,\ldots,\mathfrak p_r\}.
	\]
	为了证明每个 \(\mathfrak p_i\) 都实际出现，固定 \(i\)。不可删去 \(Q_i\) 意味着可选择
	\[
		c\in\bigcap_{j\ne i}Q_j\setminus Q_i,
	\]
	单项分解时把空交约定为 \(A\)。于是 \(c+I\) 非零，且
	\[
		\operatorname{Ann}_A(c+I)=I:c=Q_i:c.
	\]
	由\cref{g5:lem:primary-ideal-properties}，这个理想为 \(\mathfrak p_i\) 准素理想。因此循环子模 \(A(c+I)\cong A/(Q_i:c)\) 的关联素理想集恰为 \(\{\mathfrak p_i\}\)。子模的关联素理想属于原模，所以 \(\mathfrak p_i\in\operatorname{Ass}_A(A/I)\)，反向包含得证。

	现在设 \(\mathfrak p_i\) 极小。对每个 \(j\ne i\)，都有 \(\mathfrak p_j\nsubseteq\mathfrak p_i\)。取 \(s_j\in\mathfrak p_j\setminus\mathfrak p_i\)，则某个 \(s_j^{n_j}\in Q_j\)。它在 \(A_{\mathfrak p_i}\) 中为单位，故 \(Q_jA_{\mathfrak p_i}=A_{\mathfrak p_i}\)。有限交与局部化相容：若一个分式分别属于各扩张理想，对每一项选一个使分子落入该理想的分母，再乘起这有限多个分母，就使分子落入原来的交。因此
	\[
		IA_{\mathfrak p_i}
		=\bigcap_j Q_jA_{\mathfrak p_i}
		=Q_iA_{\mathfrak p_i}.
	\]
	最后由\cref{g5:lem:primary-ideal-properties}的收缩公式，\(Q_iA_{\mathfrak p_i}\) 收缩回 \(A\) 正是 \(Q_i\)。所给表达式仅涉及 \(I\) 与 \(\mathfrak p_i\)，所以孤立分量唯一。
\end{proof}

\begin{example}\label{g5:exm:nonunique-embedded-component}
	在任意域 \(k\) 上，令 \(R=k[x,y]\)、\(I=(x^2,xy)\)。同一理想有两份极小准素分解：
	\[
		I=(x)\cap(x^2,y)
		=(x)\cap(x^2,xy,y^2).
	\]
	这个例子也见\cite[第6节，习题6.6]{Matsumura1989}。逐一核验如下。\((x)\) 为素理想，而
	\[
		R/(x^2,y)\cong k[x]/(x^2),\qquad
		R/(x^2,xy,y^2)\cong k\oplus k\overline x\oplus k\overline y
	\]
	中，第二个商的理想 \((\overline x,\overline y)\) 平方为零。这两个商环中，常数项非零的元素都可用有限几何级数求逆；常数项为零的元素都幂零。因此它们的全部零因子均幂零，两个理想都为 \((x,y)\) 准素理想。

	再核对交。若 \(f=xg\in(x^2,y)\)，令 \(y=0\)，可知 \(g(x,0)\) 被 \(x\) 整除，所以 \(g\in(x,y)\)，于是 \(f\in(x^2,xy)\)。若 \(f=xg\in(x^2,xy,y^2)=(x,y)^2\)，则 \(f\) 没有一次项，迫使 \(g\) 的常数项为零，同样得到 \(g\in(x,y)\)。这证明两个交都包含于 \(I\)；反向包含从生成元直接可见。每份分解都不可删项：\(x\) 在第一项中却不在第二项中，而 \(y\) 或 \(y^2\) 分别给出第二项中不属于第一项的元素。

	两份分解的孤立分量均为 \((x)\)，关联素理想集均为 \(\{(x),(x,y)\}\)，但嵌入分量不同，因为 \(y\in(x^2,y)\) 而 \(y\notin(x^2,xy,y^2)\)。所以嵌入素理想由 \(I\) 确定，具体的嵌入准素分量可以随分解改变。
\end{example}

准素分解因而同时保留两种可用信息：由关联素理想定位零因子及嵌入部分，由孤立分量的局部化公式恢复各不可约分支上的方程。所有这些有限分解均建立在 Noether 条件上；证明中它分别保证零化子能够取极大元、素滤过能够终止，以及不可约理想证明中的零化子升链能够稳定。
% 纲要条目排为编号小节标题，并进入目录与书签；说明段仍保留于纲要注释中。

% !TeX root = ../../Book5.tex
% 纲要标题：### 2.5 局部检测与局部自由
% 纲要位置：几何学计划纲要.md，第 5102 行。
% 正文状态：初稿；尚未正式审核。

% 纲要内容开始
% 
% * 模消失和同态同构的检测
% * 有限生成投射模的局部刻画
% * 常秩、有限呈示及局部自由
% * 纤维维数不足以单独判定平坦的模型
% 
% ---
% 
% 纲要内容结束

\section{局部检测与局部自由}
\label{g5:sec:0205}

一个模可以在每个点附近用不同的生成元描述。要把这些描述用于整体问题，需要知道两件事：逐点为零是否足以推出整体为零，以及一点处的基能否在邻域内继续作为基。前一个问题由零化子和局部化解决；后一个问题还要控制生成元之间的关系，因此有限呈示条件十分重要。本节把这些局部信息与有限投射模联系起来，最后用一个带幂零元的基环说明只看纤维维数会丢掉什么。

\ProseParagraphBreak

以下环均为交换含幺环，模均为含幺模。用 \(\operatorname{Spec}A\)、\(D(f)\) 和 \(\kappa(\mathfrak p)\) 分别表示素谱、主开集和素理想处的剩余域，沿用\secref{g5:sec:0201}的约定。

\subsection{模消失和同态同构的检测}
\label{g5:subsec:0205:01}

若一个元素在 \(\mathfrak p\) 处变为零，就有某个不属于 \(\mathfrak p\) 的元素杀掉它。反过来，只要该元素原来不为零，其零化子就是一个真理想，可以放进某个极大理想中；在这个极大理想处，任何允许作分母的元素都不能杀掉它。

\begin{proposition}
\label{g5:prop:local-detection}
	设 \(A\) 为交换含幺环，\(M\) 为任意 \(A\) 模。下列条件等价：
	\begin{equivalences}
		\item \(M=0\)；
		\item 对每个素理想 \(\mathfrak p\subset A\)，有 \(M_{\mathfrak p}=0\)；
		\item 对每个极大理想 \(\mathfrak m\subset A\)，有 \(M_{\mathfrak m}=0\)。
	\end{equivalences}
	因此，对任意 \(A\) 线性同态 \(u:M\rightarrow N\)，它是单射、满射或同构，当且仅当对每个素理想处的局部化具有相应性质；也可以只在全部极大理想处检验。
\end{proposition}
\begin{proof}
	前两个蕴含由局部化定义及极大理想是素理想得到。反过来，若 \(m\in M\) 非零，则
	\[
		\operatorname{Ann}_A(m)=\{a\in A:am=0\}
	\]
	是真理想，包含在某个极大理想 \(\mathfrak m\) 中。如果 \(m/1=0\) 于 \(M_{\mathfrak m}\)，局部化的零元判据给出 \(s\notin\mathfrak m\) 且 \(sm=0\)，这与 \(s\in\operatorname{Ann}_A(m)\subseteq\mathfrak m\) 矛盾。所以一个非零模必有非零的极大理想局部化。

	由\cref{g5:thm:localization-exact}，局部化与核、余核相容。把模消失的结论分别用于 \(\ker u\) 和 \(\operatorname{coker}u\)，就得到单射和满射的判据；二者同时成立便是同构。
\end{proof}

这里没有有限生成假设。逐点局部化保留了每个元素被哪些分母杀掉的信息，足以检测整个模。同一方法还说明，两张同态相等可以逐点核验：对它们的差应用零映射判据即可。这些局部检测可对照\cite[第4节，定理4.6]{Matsumura1989}。

\ProseParagraphBreak

若把局部模再模去其极大理想，就得到纤维
\[
	M(\mathfrak p)\coloneqq M\otimes_A\kappa(\mathfrak p)
	\cong M_{\mathfrak p}/\mathfrak pM_{\mathfrak p}.
\]
\symidx{\(M(\mathfrak p)\)：模在 \(\mathfrak p\) 处的纤维}
对于有限生成 \(M\)，Nakayama 引理给出 \(M(\mathfrak p)=0\) 当且仅当 \(M_{\mathfrak p}=0\)，所以纤维也能检测消失，只检查极大理想处便够。这是\cite[第4节，定理4.8]{Matsumura1989}中的有限模版本。没有有限生成条件时则不同：作为 \(\mathbb Z\) 模的 \(\mathbb Q\) 非零，但对每个素数 \(p\)，有 \(\mathbb Q/p\mathbb Q=0\)。在同态的检测中，也应先检验局部模上的核与余核，不能把上述同构判据直接换成任意模的纤维同构判据。

\subsection{有限生成投射模的局部刻画}
\label{g5:subsec:0205:02}

自由模的一个有用性质是可以逐个提升基向量：若 \(E\rightarrow F\) 满射，则任意映射 \(A^n\rightarrow F\) 都能提升到 \(E\)。投射模把这一提升性质从选定的自由基中抽取出来。它允许模在整体上没有统一的基，而仍能在满射下分离出自身。

\begin{definition}
\label{g5:def:projective-module}
	设 \(A\) 为交换含幺环。若 \(A\) 模 \(P\) 满足：对每个满同态 \(q:E\rightarrow F\) 和每个同态 \(v:P\rightarrow F\)，存在同态 \(\widetilde v:P\rightarrow E\) 使 \(q\widetilde v=v\)，则称 \(P\) 为\term[toushemo]{投射模}[projective module]。
\end{definition}
\[
	\begin{tikzcd}[column sep=large,row sep=large]
		&P\arrow[d,"v"]\arrow[dl,dashed,"\widetilde v"']\\
		E\arrow[r,two heads,"q"']&F
	\end{tikzcd}
\]
图中所有箭头均为 \(A\) 线性映射。提升要求存在，通常不唯一。该定义与自由模直和项的描述可参看\cite[附录B，Projective and injective modules，第277页]{Matsumura1989}。

\begin{proposition}
\label{g5:prop:projective-dual-basis}
	设 \(A\) 为交换含幺环，\(P\) 为 \(A\) 模。下列条件等价：
	\begin{equivalences}
		\item \(P\) 有限生成且投射；
		\item 存在 \(A\) 模 \(Q\) 和非负整数 \(n\)，使 \(P\oplus Q\cong A^n\)；
		\item 存在有限多个 \(p_i\in P\) 和 \(\lambda_i\in\operatorname{Hom}_A(P,A)\)，使对所有 \(p\in P\) 都有
		\[
			p=\sum_{i=1}^n\lambda_i(p)p_i.
		\]
	\end{equivalences}
	此外，有限生成投射模必有限呈示。
\end{proposition}
\begin{proof}
	若 \(P\) 有限生成，选满同态 \(\pi:A^n\rightarrow P\)。投射性将 \(\operatorname{id}_P\) 提升为 \(s:P\rightarrow A^n\)，满足 \(\pi s=\operatorname{id}_P\)。于是
	\[
		A^n=s(P)\oplus\ker\pi,
	\]
	因为每个 \(a\in A^n\) 可写为 \(s\pi(a)+(a-s\pi(a))\)，两部分分别在这两个子模中，交为零。这证明第一条件蕴含第二条件。

	若 \(P\) 是 \(A^n\) 的直和项，给定 \(q:E\rightarrow F\) 和 \(v:P\rightarrow F\)，先用到 \(P\) 的投影把 \(v\) 扩成 \(A^n\rightarrow F\)，再逐个提升标准基的像，得到 \(A^n\rightarrow E\)。限制到 \(P\) 就是所需提升。\(P\) 也是 \(A^n\) 的商模，因而有限生成。

	在 \(\pi s=\operatorname{id}_P\) 中，令 \(p_i=\pi(e_i)\)，令 \(\lambda_i\) 为 \(s\) 的第 \(i\) 个坐标函数，就得到第三条件。反过来，由第三条件定义
	\[
		\pi:A^n\longrightarrow P,\quad e_i\longmapsto p_i,
		\qquad
		s:P\longrightarrow A^n,\quad
		p\longmapsto(\lambda_1(p),\ldots,\lambda_n(p)).
	\]
	则 \(\pi s=\operatorname{id}_P\)，重新得到直和分解。

	最后，\(Q\) 是 \(A^n\) 的商模，故有限生成。上述满同态 \(A^n\rightarrow P\) 的核同构于 \(Q\)，从而给出 \(P\) 的有限呈示。
\end{proof}

第三个条件称为\term[duioujibiaozhun]{对偶基标准}[dual basis criterion]。这里的 \(p_i\) 可以有关系，\(\lambda_i(p_j)\) 也未必在 \(i=j\) 时为 \(1\)、在 \(i\ne j\) 时为 \(0\)；所要求的是这些秩一算子的和等于恒等算子。若把 \(P\) 看成 \(A^n\) 的直和项，复合 \(s\pi:A^n\rightarrow A^n\) 的矩阵 \(E\) 满足 \(E^2=E\)，且其像就是 \(P\)。反过来，一个幂等矩阵也给出分解 \(A^n=\operatorname{im}E\oplus\operatorname{im}(1-E)\)。

\ProseParagraphBreak

局部环上的有限投射模可以选出真正的基。剩余向量空间先决定所需的基向量数；投射分裂随后使关系模也有限生成，从而能对关系再应用一次 Nakayama 引理。

\begin{proposition}
\label{g5:prop:projective-local-free}
	设 \((R,\mathfrak m)\) 为交换含幺局部环，\(P\) 为有限生成投射 \(R\) 模。则 \(P\) 是有限秩自由模，其秩为
	\[
		r=\dim_{R/\mathfrak m}(P/\mathfrak mP).
	\]
	因此，任意交换含幺环 \(A\) 上的有限生成投射模，在每个素理想处的局部化都是有限秩自由模。
\end{proposition}
\begin{proof}
	提升 \(P/\mathfrak mP\) 的一组基。由\cref{g5:cor:nakayama-generators}，这些提升给出满同态 \(\pi:R^r\rightarrow P\)。投射性使它分裂，所以 \(R^r\cong K\oplus P\)，其中 \(K=\ker\pi\) 是有限生成模。模去 \(\mathfrak m\) 后仍有直和分解
	\[
		(R/\mathfrak m)^r\cong K/\mathfrak mK\oplus P/\mathfrak mP.
	\]
	右侧第二项已为 \(r\) 维，故第一项为零。Nakayama 引理给出 \(K=0\)，于是 \(\pi\) 是同构。

	若 \(P\oplus Q\cong A^n\)，局部化得到 \(P_{\mathfrak p}\oplus Q_{\mathfrak p}\cong A_{\mathfrak p}^n\)，所以 \(P_{\mathfrak p}\) 仍有限生成且投射。再应用刚证的局部环结论即可。
\end{proof}

这里使用的是有限投射模的版本；其证明可对照\cite[第2节，定理2.5的有限模情形]{Matsumura1989}。下一小节将证明反向结论：有限呈示模若在每个点处自由，便是整体上的有限投射模。证明需要把点处的基扩展到开集，再将开集上的分裂拼成一个整体分裂。

\subsection{常秩、有限呈示及局部自由}
\label{g5:subsec:0205:03}

在一点选出的生成元可以扩展到邻域，这是\cref{g5:prop:generators-near-prime}的结论。但“生成”之外还要证明没有关系，才能得到自由基。下面利用有限呈示，把这项要求化为有限多个关系元同时消失。

\begin{lemma}
\label{g5:lem:local-basis-extension}
	设 \(A\) 为交换含幺环，\(M\) 为有限呈示 \(A\) 模，\(\mathfrak p\in\operatorname{Spec}A\)。若 \(M_{\mathfrak p}\) 是自由 \(A_{\mathfrak p}\) 模，则存在 \(f\notin\mathfrak p\) 使 \(M_f\) 是自由 \(A_f\) 模；给定的 \(M_{\mathfrak p}\) 的一组基可以在这样的邻域上扩展为一组基。
\end{lemma}
\begin{proof}
	因为 \(M\) 有限生成，\(M_{\mathfrak p}\) 的自由秩 \(r\) 有限。把给定基写成 \(\xi_i=m_i/s_i\)，其中 \(m_i\in M\)、\(s_i\notin\mathfrak p\)。乘以各自的单位 \(s_i\)，可知 \(m_i/1\) 也构成基。令 \(\varphi:A^r\rightarrow M\) 把标准基送到 \(m_i\)。由\cref{g5:prop:generators-near-prime}，存在 \(f_0\notin\mathfrak p\)，使 \(\varphi_{f_0}:A_{f_0}^r\rightarrow M_{f_0}\) 满射。

	有限呈示在局部化后仍有限呈示，因为可以逐项局部化一个有限关系矩阵。故由\cref{g5:prop:finite-presentation-kernel}，\(K=\ker\varphi_{f_0}\) 是有限生成 \(A_{f_0}\) 模。它在 \(\mathfrak pA_{f_0}\) 处的局部化为零，因为 \(\varphi_{\mathfrak p}\) 是基给出的同构。逐个杀掉 \(K\) 的有限生成元，再把所用分母相乘，得到 \(a\notin\mathfrak p\)，使 \(K\) 在进一步反演 \(a\) 后为零。具体说，这些分母在 \(A_{f_0}\) 中可写成 \(a_j/f_0^{e_j}\)，其中 \(a_j\notin\mathfrak p\)，取 \(a=\prod_j a_j\) 即可。由局部化正合性，\(\varphi_{f_0a}\) 是同构。

	最后令 \(f=f_0a\prod_i s_i\)。所有 \(s_i\) 在 \(A_f\) 中为单位，所以再将基 \(m_i/1\) 分别除以 \(s_i\)，就得到扩展原基 \(\xi_i\) 的自由基。
\end{proof}

这同时证明有限呈示模的自由点组成开集，参看\cite[第4节，定理4.10(ii)]{Matsumura1989}。证明中有限生成用于控制余核，有限呈示用于控制选定满射的核；两个有限性条件对应不同的步骤。

\ProseParagraphBreak

若模可以在覆盖素谱的主开集上分别选出有限自由基，就称其\term[jubuziyou]{局部自由}[locally free]。在本章的模语言中，这具体指一个主开覆盖 \(D(f_i)\) 和同构 \(M_{f_i}\cong A_{f_i}^{r_i}\)；秩 \(r_i\) 可以随开集改变。

\begin{theorem}
\label{g5:thm:finite-projective-local}
	设 \(A\) 为交换含幺环，\(M\) 为有限呈示 \(A\) 模。下列条件等价：
	\begin{equivalences}
		\item \(M\) 是有限生成投射模；
		\item 对每个素理想 \(\mathfrak p\)，\(M_{\mathfrak p}\) 是自由 \(A_{\mathfrak p}\) 模；
		\item 对每个极大理想 \(\mathfrak m\)，\(M_{\mathfrak m}\) 是自由 \(A_{\mathfrak m}\) 模；
		\item 存在有限主开覆盖 \(\operatorname{Spec}A=\bigcup_{i=1}^tD(f_i)\)，使每个 \(M_{f_i}\) 是有限秩自由 \(A_{f_i}\) 模。
	\end{equivalences}
\end{theorem}
\begin{proof}
	第一条件蕴含第二条件，已由\cref{g5:prop:projective-local-free}证明；第二条件自然蕴含第三条件。若第三条件成立，由\cref{g5:lem:local-basis-extension}，每个极大理想都有一个主开邻域，使 \(M\) 在其中自由。这些主开集的定义元素生成单位理想，否则某个极大理想会包含它们全部，与其已被覆盖矛盾。于是 \(1\) 是有限多个定义元素的线性组合，相应的有限多个主开集就覆盖整个素谱，得到第四条件。

	最后设第四条件成立，选满同态 \(\pi:F=A^n\rightarrow M\)。因为 \(M_{f_i}\) 自由，可提升其自由基，得到一个局部分裂
	\[
		s_i:M_{f_i}\longrightarrow F_{f_i},\qquad
		\pi_{f_i}s_i=\operatorname{id}_{M_{f_i}}.
	\]
	我们要把它变成整体映射，再组合各个局部分裂。由 \(M\) 有限呈示及\cref{g5:prop:hom-localization}，
	\[
		\operatorname{Hom}_A(M,F)_{f_i}
		\cong\operatorname{Hom}_{A_{f_i}}(M_{f_i},F_{f_i}),
	\]
	所以可以写 \(s_i=\tau_i/f_i^{a_i}\)，其中 \(\tau_i:M\rightarrow F\) 为整体同态。等式 \(\pi_{f_i}s_i=\operatorname{id}\) 表明
	\(\pi\tau_i-f_i^{a_i}\operatorname{id}_M\) 在反演 \(f_i\) 后为零。再对 \(\operatorname{Hom}_A(M,M)\) 使用同一局部化结论，找到 \(b_i\geq0\) 使
	\[
		\pi(f_i^{b_i}\tau_i)=f_i^{a_i+b_i}\operatorname{id}_M
		\quad\text{在整个 }M\text{ 上成立}.
	\]
	可再乘一次 \(f_i\)，使指数 \(N_i\) 为正。记相应调整后的同态为 \(\rho_i\)，于是 \(\pi\rho_i=f_i^{N_i}\operatorname{id}_M\)。

	由于 \(D(f_i^{N_i})=D(f_i)\) 仍覆盖素谱，存在 \(c_i\in A\) 使 \(\sum_i c_i f_i^{N_i}=1\)。因此
	\[
		s=\sum_i c_i\rho_i:M\longrightarrow F,\qquad
		\pi s=\sum_i c_i f_i^{N_i}\operatorname{id}_M=\operatorname{id}_M.
	\]
	这给出整体分裂。由\cref{g5:prop:projective-dual-basis}，\(M\) 是有限生成投射模。
\end{proof}

局部基不必在交集上逐项相同；上面拼合的是满射的分裂。有限呈示让局部同态可以清除分母成为整体同态，而主开覆盖让相应分母的幂产生单位理想。这一局部刻画也见于\cite[第7节，定理7.12]{Matsumura1989}。

\ProseParagraphBreak

在自由邻域上，纤维的维数就是自由基的项数。因此，有限投射模具有一个逐点可读的\term[zhihanshu]{秩函数}[rank function]：
\[
	\operatorname{rk}_M(\mathfrak p)
	\coloneqq\dim_{\kappa(\mathfrak p)}M(\mathfrak p).
\]
\symidx{\(\operatorname{rk}_M(\mathfrak p)\)：有限投射模在一点的秩}
若 \(M_f\cong A_f^r\)，则这个函数在 \(D(f)\) 上恒等于 \(r\)，所以秩函数局部常值。其每个取值的原像都是开集，补集也是其他取值的原像之并，因而也是闭集。特别地，若 \(\operatorname{Spec}A\) 非空且连通，秩在全空间恒定。

\begin{example}
\label{g5:exm:projective-varying-rank}
	设 \(k\) 为域，\(A=k\times k\)，并取幂等元 \(e=(1,0)\)。模 \(P=Ae\) 满足
	\[
		A=Ae\oplus A(1-e),
	\]
	所以它有限生成且投射。\(A\) 的两个素理想是 \(\mathfrak p_1=0\times k\) 和 \(\mathfrak p_2=k\times0\)：素理想必须包含 \(e,1-e\) 之一，因为它们的乘积为零，而不能同时包含二者。模 \(P\) 在第一点处为一维自由模，在第二点处为零模，秩分别是 \(1,0\)。这两个点都是开闭集，故秩函数仍局部常值，但不存在整数 \(r\) 使 \(P\cong A^r\)。
\end{example}

对于一般有限呈示模，纤维维数可能在特殊点增大。一个关系矩阵可以把这种变化的发生位置写成子式方程。

\begin{proposition}
\label{g5:prop:fiber-upper-semicontinuity}
	设 \(A\) 为交换含幺环，\(M\) 有有限呈示
	\[
		A^s\xrightarrow{H}A^n\longrightarrow M\longrightarrow0.
	\]
	记 \(H\) 的 \(q\) 阶子式生成的理想为 \(J_q(H)\)，不存在这种子式时规定 \(J_q(H)=0\)。则
	\[
		\nu_M(\mathfrak p)\coloneqq
		\dim_{\kappa(\mathfrak p)}M(\mathfrak p)
		=n-\operatorname{rank}_{\kappa(\mathfrak p)}H(\mathfrak p),
	\]
	且对 \(1\leq r\leq n\)，有
	\[
		\{\mathfrak p:\nu_M(\mathfrak p)\geq r\}
		=V(J_{n-r+1}(H)).
	\]
	因此纤维维数函数\term[shangbanlianxu]{上半连续}[upper semicontinuous]，即它大于等于给定整数的点集为闭集。
\end{proposition}
\begin{proof}
	由\cref{g5:prop:tensor-right-exact}，将呈示张量 \(\kappa(\mathfrak p)\) 后得到
	\[
		\kappa(\mathfrak p)^s\xrightarrow{H(\mathfrak p)}
		\kappa(\mathfrak p)^n\longrightarrow M(\mathfrak p)\longrightarrow0.
	\]
	向量空间的秩公式给出所述维数。它至少为 \(r\)，当且仅当矩阵的秩至多为 \(n-r\)，也就是所有 \(n-r+1\) 阶子式在 \(\kappa(\mathfrak p)\) 中为零。一个 \(A\) 中元素在 \(\kappa(\mathfrak p)\) 中为零，当且仅当它属于 \(\mathfrak p\)。所以这一条件恰是 \(J_{n-r+1}(H)\subseteq\mathfrak p\)。当 \(r\leq0\) 或 \(r>n\) 时，相应集合分别为全空间或空集，仍为闭集。
\end{proof}

例如 \(A=k[t]\)、\(M=A/(t)\) 有关系矩阵 \((t)\)。在 \(t\ne0\) 的点处，纤维为零；在 \((t)\) 处，纤维维数为 \(1\)。关系的系数在该点消失，使纤维中可见的关系减少，因此维数增大。有限生成模的纤维维数也上半连续：在一点选出至多 \(r\) 个生成元后，\cref{g5:prop:generators-near-prime}使它们在邻域内继续生成；结合 Nakayama 引理，便知 \(\{\nu_M\leq r\}\) 为开集。有限呈示则进一步给出上面的有限矩阵与子式公式。

\subsection{纤维维数不足以单独判定平坦的模型}
\label{g5:subsec:0205:04}

局部自由模在每个邻域都像固定数目的独立坐标，因此与它作张量积会保留单射。平坦性抽取的正是这个性质：变换系数时，原来嵌在另一个模中的关系不能突然被压成零。\secref{g5:sec:0202}已经说明张量积总是右正合，因而还需要控制的是左端的单射。

\begin{definition}
\label{g5:def:flat-module}
	设 \(A\) 为交换含幺环。若 \(A\) 模 \(M\) 满足：对每个单同态 \(L\hookrightarrow N\)，诱导映射
	\[
		L\otimes_A M\longrightarrow N\otimes_A M
	\]
	仍为单射，则称 \(M\) 为\term[pingtanmo]{平坦模}[flat module]。等价地，函子 \(-\otimes_A M\) 保持短正合列。
\end{definition}

自由模 \(A^r\) 平坦，因为张量后只是取 \(r\) 份原映射的直和。有限投射模也平坦：若 \(P\oplus Q\cong A^r\)，那么与 \(P\) 张量后的映射是与 \(A^r\) 张量后所得单射的一个直和分量。局部化 \(A_f\) 作为 \(A\) 模同样平坦，这是张量积与局部化相容以及局部化正合性的直接结果。

\begin{example}
\label{g5:exm:constant-fiber-not-flat}
	设 \(k\) 为域，\(A=k[\varepsilon]/(\varepsilon^2)\)，\(M=A/(\varepsilon)\)。素谱只有 \(\mathfrak m=(\varepsilon)\) 一个点，剩余域为 \(k\)，且
	\[
		M\otimes_A k\cong k.
	\]
	所以纤维维数在整个基上恒为 \(1\)。然而 \(M\) 不平坦。取理想 \(I=(\varepsilon)\) 的包含 \(I\hookrightarrow A\)，张量 \(M\) 后得到
	\[
		I\otimes_A M\cong I/\varepsilon I=I
		\longrightarrow A\otimes_A M\cong M.
	\]
	左端非零，映射却把 \(\varepsilon\otimes\overline1\) 送到 \(\varepsilon\overline1=0\)，因而不是单射。纤维只看到剩余域，丢掉了基环中的非零幂零方向；平坦性同时检验这个方向与模之间的关系。
\end{example}

这个模甚至有限呈示：\(A\xrightarrow{\varepsilon}A\rightarrow M\rightarrow0\) 就是一个有限呈示。因此在允许非约化基环时，有限呈示与纤维维数局部常值合在一起，仍不能替代平坦性。要理解平坦与局部自由的联系，应直接分析张量积对关系模的作用。

\begin{proposition}
\label{g5:prop:flat-local-detection}
	设 \(A\) 为交换含幺环，\(M\) 为任意 \(A\) 模。则 \(M\) 平坦，当且仅当每个 \(M_{\mathfrak p}\) 都是平坦 \(A_{\mathfrak p}\) 模。也可以只检查全部极大理想。
\end{proposition}
\begin{proof}
	若 \(M\) 平坦，取 \(A_{\mathfrak p}\) 模的单射 \(L\hookrightarrow N\)，将其视为 \(A\) 模的单射。对任意 \(A_{\mathfrak p}\) 模 \(L\)，有自然同构
	\[
		L\otimes_A M\cong L\otimes_{A_{\mathfrak p}}M_{\mathfrak p},
	\]
	因为 \(A\setminus\mathfrak p\) 中的元素已经在 \(L\) 上可逆，纯张量的分母可以移到 \(L\) 一侧。于是 \(M\) 的平坦性给出 \(M_{\mathfrak p}\) 的平坦性。

	反过来，设所有极大理想处的局部模都平坦。对任意单射 \(L\hookrightarrow N\)，记诱导映射 \(L\otimes_A M\rightarrow N\otimes_A M\) 的核为 \(K\)。将它在极大理想 \(\mathfrak m\) 处局部化，利用局部化的正合性及其与张量积相容，得到
	\[
		K_{\mathfrak m}
		=\ker\bigl(L_{\mathfrak m}\otimes_{A_{\mathfrak m}}M_{\mathfrak m}
		\longrightarrow N_{\mathfrak m}\otimes_{A_{\mathfrak m}}M_{\mathfrak m}\bigr)=0.
	\]
	最后一个等号来自局部化仍保持原单射，以及 \(M_{\mathfrak m}\) 平坦。由\cref{g5:prop:local-detection}，\(K=0\)，故 \(M\) 平坦。
\end{proof}

下面的证明用一个具体等式控制关系。若 \(0\rightarrow K\rightarrow F\rightarrow M\rightarrow0\) 正合，平坦性保证 \(K\) 与 \(IF\) 的交没有比 \(IK\) 更多的元素。取局部环的极大理想作为 \(I\)，就可以把“极小生成元”提升为“自由基”。

\begin{theorem}
\label{g5:thm:finite-presented-flat}
	设 \(A\) 为交换含幺环，\(M\) 为有限呈示 \(A\) 模。则以下条件等价：\(M\) 平坦；\(M\) 局部自由；\(M\) 是有限生成投射模。
\end{theorem}
\begin{proof}
	由\cref{g5:thm:finite-projective-local}，后两个条件等价；有限投射模平坦已经在定义后证明。因此只需证明，有限呈示平坦模在每个素理想处自由。

	先设 \(M\) 为任意环上的平坦模，且 \(0\rightarrow K\rightarrow F\xrightarrow{\pi}M\rightarrow0\) 正合。对任意理想 \(I\subseteq A\)，我们证明
	\begin{equation}
	\label{g5:eq:flat-kernel-intersection}
		K\cap IF=IK.
	\end{equation}
	包含 \(IK\subseteq K\cap IF\) 直接成立。若 \(x\in K\cap IF\)，写 \(x=\sum_j a_jv_j\)，其中 \(a_j\in I\)、\(v_j\in F\)。张量
	\(z=\sum_j a_j\otimes v_j\in I\otimes_A F\)
	在 \(I\otimes_A M\) 中的像，经乘法映射到 \(M\) 后为 \(\pi(x)=0\)。平坦性使 \(I\otimes_A M\rightarrow A\otimes_A M=M\) 为单射，所以 \(z\) 在 \(I\otimes_A M\) 中的像已经为零。再用张量积的右正合性，
	\[
		I\otimes_A K\longrightarrow I\otimes_A F
		\longrightarrow I\otimes_A M\longrightarrow0,
	\]
	可知 \(z\) 来自 \(I\otimes_A K\)。对这些张量作标量乘法，就得到 \(x\in IK\)，证明了\eqref{g5:eq:flat-kernel-intersection}。

	现在固定 \(\mathfrak p\)，令 \(R=A_{\mathfrak p}\)、\(\mathfrak m=\mathfrak pA_{\mathfrak p}\)、\(P=M_{\mathfrak p}\)。模 \(P\) 仍有限呈示，并由\cref{g5:prop:flat-local-detection}知它平坦。提升 \(P/\mathfrak mP\) 的一组基，得到满同态 \(R^r\rightarrow P\)。记其核为 \(K\)。由\cref{g5:prop:finite-presentation-kernel}，\(K\) 有限生成；由\cref{g5:prop:minimal-presentation-local}，\(K\subseteq\mathfrak mR^r\)。把\eqref{g5:eq:flat-kernel-intersection}用于 \(I=\mathfrak m\)、\(F=R^r\)，便得到
	\[
		K=K\cap\mathfrak mR^r=\mathfrak mK.
	\]
	Nakayama 引理使 \(K=0\)，因此 \(P\cong R^r\)。这对所有素理想成立，再由\cref{g5:thm:finite-projective-local}得到结论。
\end{proof}

平坦性的局部检测与有限呈示模的上述等价关系可对照\cite[第7节，定理7.1、7.10、7.12及其推论]{Matsumura1989}。在\cref{g5:exm:constant-fiber-not-flat}中，极小满射 \(A\rightarrow M\) 的核为 \(K=(\varepsilon)\)，虽然 \(K\subseteq\mathfrak mA\)，却有 \(\mathfrak mK=0\ne K\)。这恰好指出平坦证明在哪里失败：纤维中看不见的关系 \(\varepsilon\overline1=0\) 仍保留在局部模中。

% 纲要条目排为编号小节标题，并进入目录与书签；说明段仍保留于纲要注释中。


\begin{chaptersummary}
	局部化由可用分母决定，其等价关系保留零因子造成的消失现象。素理想处的局部化是局部环，剩余域把该点的函数值组织成一个域；理想的扩缩则把局部化与饱和联系起来。素谱上的主开集为这些运算提供了邻域语言。

	\ProseParagraphBreak

	有限生成模只要求有限个生成方向，有限呈示还要求有限条生成关系。局部化保持正合，张量积一般只保持右正合，Hom 一般只保持左正合。有限呈示使 Hom 与局部化相容，因而允许把局部同态表示为有共同分母的整体同态。

	\ProseParagraphBreak

	Nakayama 引理把有限模的局部生成化为剩余域上的线性代数。支集说明模在哪些素理想处非零，关联素理想通过元素的零化子识别零因子和嵌入分支；Noether 环上的有限模只有有限个关联素理想，理想则能分解为有限个准素理想的交。

	\ProseParagraphBreak

	模是否为零、同态是否为同构，都可以在所有极大理想处检测。对于有限呈示模，逐点自由可以扩展成有限个主开集上的自由表示，并拼出整体投射结构。秩在局部自由情形下局部常值；仅有恒定的纤维维数仍可能遗漏幂零方向中的关系，平坦性正要求张量积保留原有的单射。
\end{chaptersummary}

\BookExercises

\begin{exercise}\label{g5:ex:localize-twelve}
	设 \(A=\mathbb Z/12\mathbb Z\)。分别计算 \(A_{\overline2}\) 和 \(A_{\overline6}\)，写出自然同态的核，并说明在第一个环中为何 \(\overline3\) 变为零而 \(\overline4\) 没有变为零。
\end{exercise}
\begin{solution}
	模三同态把 \(\overline2\) 送到单位，因而诱导 \(A_{\overline2}\to\mathbb F_3\)。在 \(A_{\overline2}\) 中， \(4\cdot3=0\)，且 \(4\) 为单位，所以 \(3=0\)。 因而自然同态经 \(\mathbb F_3\) 分解，两个方向互逆，得到 \(A_{\overline2}\cong\mathbb F_3\)，其核为 \((\overline3)\)。 因此 \(\overline3\) 变为零，而 \(\overline4\) 对应 \(\overline1\in\mathbb F_3\)，没有变为零。 直接用消失判据也可检查：\(4\cdot3=0\)，但 \(2^n4\) 从不被 \(12\) 整除。

	又因 \(\overline6^{\,2}=0\)，允许的分母中含零， \(A_{\overline6}=0\)，自然同态的核为整个 \(A\)。
\end{solution}

\begin{exercise}\label{g5:ex:localization-finite-intersection}
	设 \(A\) 为交换含幺环，\(S\subseteq A\) 为乘法集，\(N_1,N_2\) 为 \(A\)-模 \(M\) 的子模。证明
	\[
	S^{-1}(N_1\cap N_2)
	=(S^{-1}N_1)\cap(S^{-1}N_2)
	\]
	于 \(S^{-1}M\) 中成立。 再以 \(A=\mathbb Z\)、\(S=\{2^n:n\geq0\}\)、\(N_j=2^j\mathbb Z\) 说明这个等式不能直接推广为无限交。
\end{exercise}
\begin{solution}
	左侧包含于右侧。若 \(m/s\) 属于右侧，分式等价关系给出 \(u,v\in S\) 使 \(um\in N_1\)、\(vm\in N_2\)。 于是 \(uvm\in N_1\cap N_2\)，且 \(m/s=uvm/(uvs)\)，得到反向包含。 对给定无限族，\(\bigcap_{j\geq0}2^j\mathbb Z=0\)，所以先取交再局部化得零； 但每个 \(S^{-1}N_j\) 都等于 \(\mathbb Z[1/2]\)，局部化后再取交仍为该非零环。
\end{solution}

\begin{exercise}\label{g5:ex:tensor-hom-cyclic}
	设 \(m,n\) 为正整数。计算
	\[
	(\mathbb Z/m\mathbb Z)\otimes_{\mathbb Z}(\mathbb Z/n\mathbb Z),
	\qquad
	\operatorname{Hom}_{\mathbb Z}(\mathbb Z/m\mathbb Z,\mathbb Z/n\mathbb Z),
	\]
	并说明两者虽有同构的抽象群，所表示的数据却不同。
\end{exercise}
\begin{solution}
	令 \(d=\gcd(m,n)\)。由 \((\mathbb Z/m)\otimes N\cong N/mN\)，张量积为 \(\mathbb Z/(m,n)\cong\mathbb Z/d\mathbb Z\)，由 \(\overline1\otimes\overline1\) 生成，记录双线性配对的泛对象。

	一个同态由 \(\overline1\) 的像 \(b\in\mathbb Z/n\mathbb Z\) 唯一决定，条件是 \(mb=0\)。 写 \(m=dm'\)、\(n=dn'\)，其中 \(\gcd(m',n')=1\)，则 \(n\mid mb\) 等价于 \(n'\mid b\)。 因而 Hom 群由 \(n'\) 的剩余类生成，共有 \(d\) 个元素，也同构于 \(\mathbb Z/d\mathbb Z\)。 前者是商模，后者是被 \(m\) 杀掉的子模；两者在自然性与正合性中的作用仍然不同。
\end{solution}

\begin{exercise}\label{g5:ex:ideal-two-variable-presentation}
	设 \(k\) 为域，\(A=k[x,y]\)，\(I=(x^2,xy,y^2)\)。证明下列序列正合：
	\[
	0\longrightarrow A^2
	\xrightarrow{\ D\ } A^3
	\xrightarrow{\ (x^2\;\;xy\;\;y^2)\ } I
	\longrightarrow0,\qquad
	D=\begin{pmatrix}-y&0\\x&-y\\0&x\end{pmatrix}.
	\]
	并计算 \(I\otimes_A A/(x,y)\) 的维数，比较它与 \(I\otimes_A k(x,y)\) 的维数。
\end{exercise}
\begin{solution}
	右侧满性来自 \(I\) 的定义，两个箭头复合为零。 若 \(ax^2+bxy+cy^2=0\)，令 \(x=0\)，得 \(c(0,y)y^2=0\)，所以 \(c=x\beta\)。 在整环中约去非零 \(x\)，得到 \(ax+by+\beta y^2=0\)； 再令 \(y=0\)，可知 \(a=-y\alpha\) 对某个 \(\alpha\in A\) 成立。 代回并约去 \(y\)，得到 \(b=x\alpha-y\beta\)。 因而全部关系都由 \(D\) 的两列生成。 若 \(D(\alpha,\beta)=0\)，第一和第三个坐标分别给出 \(y\alpha=0\)、\(x\beta=0\)，故 \(\alpha=\beta=0\)，左侧单射。

	张量积 \(I\otimes_A A/(x,y)\cong I/(x,y)I\) 的基是 \(x^2,xy,y^2\) 的像，维数为三；其独立性由比较二次齐次项得到。 在 \(k(x,y)\) 中 \(x^2\) 为单位，故扩张理想等于整个域，维数为一。 生成方向的数量在原点增加，说明有限呈示模的纤维维数可以跳变。
\end{solution}

\begin{exercise}\label{g5:ex:hom-nonsplit-cyclic}
	设整数 \(n\geq2\)。证明短正合列
	\[
	0\longrightarrow\mathbb Z\xrightarrow{\ \times n\ }\mathbb Z
	\longrightarrow\mathbb Z/n\mathbb Z\longrightarrow0
	\]
	不分裂。分别用 Hom 和张量积指出其中一个满射或单射在相应运算后丢失。
\end{exercise}
\begin{solution}
	若有截面 \(\mathbb Z/n\to\mathbb Z\)，\(\overline1\) 的像被 \(n\) 杀掉，故只能为零，不能映回 \(\overline1\)，所以不分裂。 应用 \(\operatorname{Hom}_{\mathbb Z}(\mathbb Z/n,-)\) 后， \(\operatorname{Hom}(\mathbb Z/n,\mathbb Z)=0\)，而 \(\operatorname{Hom}(\mathbb Z/n,\mathbb Z/n)\) 含非零恒等同态； 因此最后的箭头不满。 应用 \(-\otimes_{\mathbb Z}\mathbb Z/n\) 后，原来的乘 \(n\) 单射变成 \(\mathbb Z/n\xrightarrow{0}\mathbb Z/n\)，不再单射。
\end{solution}

\begin{exercise}\label{g5:ex:surjective-finite-endomorphism}
	设 \(A\) 为交换含幺环，\(M\) 为有限生成 \(A\)-模。 证明每个满的 \(A\)-线性同态 \(\varphi:M\to M\) 都是同构。 提示：令多项式变量 \(T\) 通过 \(\varphi\) 作用于 \(M\)，对 \(A[T]\) 和理想 \((T)\) 使用\cref{g5:prop:nakayama-ideal}。
\end{exercise}
\begin{solution}
	赋予 \(M\) 以 \(A[T]\)-模结构 \(f(T)m=f(\varphi)m\)。原来的有限个 \(A\)-生成元也生成这个 \(A[T]\)-模。 满性给出 \(TM=M\)，故存在 \(Tq(T)\in(T)\) 使 \((1-Tq(T))M=0\)。 换成算子就是 \(\operatorname{id}_M=\varphi q(\varphi)\)。 多项式 \(q(\varphi)\) 与 \(\varphi\) 交换，因此也是左逆； \(\varphi\) 为同构，逆为 \(q(\varphi)\)。 这个论证不要求 \(A\) Noether。
\end{solution}

\begin{exercise}\label{g5:ex:generators-after-residue}
	设 \((A,\mathfrak m)\) 为交换含幺局部环，\(k=A/\mathfrak m\)， \(u:M\to N\) 为有限生成 \(A\)-模之间的同态。 证明 \(u\) 满当且仅当 \(u\otimes_A k:M\otimes_Ak\to N\otimes_Ak\) 满。 再举例说明：剩余域上的同态为同构，不能保证 \(u\) 为同构。
\end{exercise}
\begin{solution}
	令 \(Q=\operatorname{coker}u\)。张量积的右正合性给出 \(\operatorname{coker}(u\otimes_Ak)\cong Q/\mathfrak mQ\)。 \(Q\) 为有限生成模，Nakayama 引理说明该商为零当且仅当 \(Q=0\)，即得满性判据。 取 \(A=k[t]_{(t)}\)、\(M=A\)、\(N=A/(t)\)，令 \(u\) 为商映射。 剩余域上的同态是 \(k\) 的恒等同态，然而 \(u\) 的核为非零理想 \((t)\)。
\end{solution}

\begin{exercise}\label{g5:ex:support-without-finiteness}
	作为 \(\mathbb Z\)-模，令 \(M=\bigoplus_{p\text{ 为素数}}\mathbb Z/p\mathbb Z\)。 求 \(\operatorname{Ann}_{\mathbb Z}M\) 和 \(\operatorname{Supp}_{\mathbb Z}M\)，证明其支集不是闭集，并指出有限模支集公式在这里为何不能使用。
\end{exercise}
\begin{solution}
	杀掉整个 \(M\) 的整数必须被所有素数整除，所以 \(\operatorname{Ann}M=0\)。 在 \((0)\) 处局部化时，每个元素仅有有限个非零分量，故被某个非零整数杀掉；因此 \(M_{(0)}=0\)。 在 \((p)\) 处，第 \(p\) 个分量仍为 \(\mathbb Z/p\)，而其余分量被可逆素数杀掉，故 \(M_{(p)}\cong\mathbb Z/p\)。 支集于是是所有非零素理想 \((p)\)，不含 \((0)\)。

	一个闭集若含所有 \((p)\)，其定义理想便包含于 \(\bigcap_p(p)=0\)，所以它只能是整个 \(\operatorname{Spec}\mathbb Z\)。 因而这个支集不闭，也不等于 \(V(\operatorname{Ann}M)=\operatorname{Spec}\mathbb Z\)。 \(M\) 不有限生成，因为任何有限个元素只能涉及有限多个直和分量；缺少有限性时，不能把逐个元素的消失分母统一起来。
\end{solution}

\begin{exercise}\label{g5:ex:primary-product-powers}
	设 \(k\) 为域，\(A=k[x,y]\)，\(I=(x^2y)\)。 证明
	\[
	I=(x^2)\cap(y)
	\]
	是不可删减的准素分解，求 \(\operatorname{Ass}_A(A/I)\)，并判断是否有嵌入素理想。
\end{exercise}
\begin{solution}
	若 \(x^2a\in(y)\)，令 \(y=0\) 得 \(x^2a(x,0)=0\)，所以 \(y\mid a\)，从而 \((x^2)\cap(y)=(x^2y)\)。 理想 \((y)\) 素。 对 \((x^2)\)，把非零多项式写成 \(x^rf\)、\(x\nmid f\)； 关于 \(x\) 的最低次数在乘法下相加，因为系数环 \(k[y]\) 是整环。 因而若 \(ab\in(x^2)\) 且 \(a\notin(x)\)，则 \(b\in(x^2)\)，这说明 \((x^2)\) 是 \((x)\)-准素的。

	两个分量都不可删，根式不同，故此分解最简。 由准素分解的关联素理想结论， \(\operatorname{Ass}_A(A/I)=\{(x),(y)\}\)。 也可直接看到 \(\operatorname{Ann}(\overline{xy})=(x)\)、 \(\operatorname{Ann}(\overline{x^2})=(y)\)。 两个素理想都极小地包含 \(I\)，所以没有嵌入素理想；\(x^2\) 的幂次记录沿一个分支的非约化结构，本身不产生嵌入分支。
\end{solution}

\begin{exercise}\label{g5:ex:all-fibres-zero-infinite}
	证明非零 \(\mathbb Z\)-模 \(M=\mathbb Q/\mathbb Z\) 满足
	\[
	M\otimes_{\mathbb Z}\kappa(\mathfrak p)=0
	\quad\text{对每个 }\mathfrak p\in\operatorname{Spec}\mathbb Z.
	\]
	再证明 \(M_{(p)}\ne0\) 对每个素数 \(p\) 成立，以区分纤维与局部化。
\end{exercise}
\begin{solution}
	零素理想的剩余域是 \(\mathbb Q\)。每个 \(M\) 的元素都被某个非零整数杀掉，因此 \(M\otimes\mathbb Q=0\)。 对 \((p)\)，剩余域是 \(\mathbb F_p\)，且 \(M\otimes\mathbb F_p\cong M/pM=0\)，因为任何有理数的剩余类都可除以 \(p\)。 但 \(1/p+\mathbb Z\) 在 \(M_{(p)}\) 中非零：若被某个 \(s\notin(p)\) 杀掉，就会有 \(s/p\in\mathbb Z\)，即 \(p\mid s\)，矛盾。 有限模的剩余域检测来自 Nakayama 引理，这个非有限生成模不满足其条件；对任意模成立的局部化检测仍然有效。
\end{solution}

\begin{exercise}\label{g5:ex:module-two-open-gluing}
	设 \(A\) 为交换含幺环，\(f,g\in A\) 满足 \((f,g)=A\)，\(M\) 为任意 \(A\)-模。 证明序列
	\[
	0\longrightarrow M\longrightarrow M_f\oplus M_g
	\xrightarrow{\ (u,v)\mapsto u|_{fg}-v|_{fg}\ }M_{fg}
	\longrightarrow0
	\]
	正合，其中第一个非零箭头是同时局部化。
\end{exercise}
\begin{solution}
	先注意 \((f^r,g^s)=A\) 对正整数 \(r,s\) 成立：若这个理想真包含于某个极大理想，素性便迫使该极大理想同时含 \(f,g\)，矛盾。 若 \(m\) 在 \(M_f,M_g\) 中都为零，存在 \(r,s\) 使 \(f^rm=g^sm=0\)；用 \(cf^r+dg^s=1\) 得 \(m=0\)，所以第一个箭头单射。

	若 \(u=m/f^r\)、\(v=n/g^s\) 在 \(M_{fg}\) 中相同，则某个 \(t\geq0\) 满足 \((fg)^t(g^sm-f^rn)=0\)。 取充分大 \(N\)，令 \(a=f^{N-r}m\)、\(b=g^{N-s}n\)，便有 \(u=a/f^N\)、\(v=b/g^N\) 且 \(g^Na=f^Nb\) 于 \(M\) 中成立。 选 \(c,d\) 使 \(cf^N+dg^N=1\)，置 \(w=ca+db\)。 在 \(M_f\) 中 \(b=g^Nu\)、\(a=f^Nu\)，所以 \(w=u\)；在 \(M_g\) 中同理 \(w=v\)。 这证明中间正合。

	任取 \(m/(fg)^N\in M_{fg}\)，仍选 \(cf^N+dg^N=1\)，有
	\[
	\frac{m}{f^Ng^N}=\frac{dm}{f^N}+\frac{cm}{g^N}.
	\]
	它是 \((dm/f^N,-cm/g^N)\) 的像，故最后一个箭头满。
\end{solution}

\begin{exercise}\label{g5:ex:idempotent-variable-rank}
	设 \(k\) 为域，\(A=k\times k\times k\)、\(e=(1,0,1)\)，\(P=Ae\)。 证明 \(P\) 是有限生成投射模，求其各点的秩，并证明它不是自由模。
\end{exercise}
\begin{solution}
	\(e^2=e\)，且 \(A=Ae\oplus A(1-e)\)，所以 \(P\) 是秩一有限自由模 \(A\) 的直和项，为有限生成投射模。 环 \(A\) 的素理想恰为三个坐标投影的核 \(\mathfrak m_1,\mathfrak m_2,\mathfrak m_3\)。 事实上三个坐标幂等元两两乘积为零、和为 \(1\)，在整环 \(A/\mathfrak p\) 中恰有一个像为 \(1\)，其余为零，从而给出上述三个核。 在三个点的局部化分别等于 \(k\)，而 \(e\) 的像分别为 \(1,0,1\)，故 \(P\) 的秩分别为 \(1,0,1\)。 每个点都是开点，所以秩局部常值；若 \(P\cong A^r\)，三个点的秩都应等于同一个 \(r\)，矛盾。
\end{solution}

\begin{exercise}\label{g5:ex:projective-cyclic-idempotent}
	设 \(A\) 为交换含幺环，\(I\subseteq A\) 为理想。 证明 \(A/I\) 为投射 \(A\)-模，当且仅当 \(I=Ae\) 对某个幂等元 \(e^2=e\) 成立。
\end{exercise}
\begin{solution}
	若 \(A/I\) 投射，商映射有截面，因而 \(A=I\oplus J\) 为 \(A\)-模的直和。令 \(\pi:A\to I\) 为对应投影，\(e=\pi(1)\)。 \(A\)-线性给出 \(\pi(a)=ae\)，故 \(I=Ae\)； 又因 \(\pi\) 在 \(I\) 上恒等，\(\pi(e)=e\)，即 \(e^2=e\)。

	反过来，若 \(e^2=e\)，则 \(A=Ae\oplus A(1-e)\)：两项之和含 \(1\)，且若 \(z=ae=b(1-e)\)，乘以 \(e\) 得 \(z=0\)。 因此 \(A/I\cong A(1-e)\) 是有限自由模 \(A\) 的直和项，为投射模。
\end{solution}

\begin{exercise}\label{g5:ex:reduced-constant-fibre}
	设 \(A\) 为约化交换含幺环，\(M\) 为有限呈示 \(A\)-模。 若函数 \(\mathfrak p\mapsto\dim_{\kappa(\mathfrak p)} (M\otimes_A\kappa(\mathfrak p))\) 局部常值，证明 \(M\) 局部自由。 说明本题与双数环上的恒定纤维维数反例为何不矛盾。
\end{exercise}
\begin{solution}
	固定 \(\mathfrak p\)，在一个主开邻域上令纤维维数恒为 \(r\)，并把环替换为这个主开集对应的局部化。 有限呈示写成 \(A^q\xrightarrow C A^n\to M\to0\)。 对每个素理想，矩阵 \(C\) 在剩余域上的秩都是 \(s=n-r\)。 在 \(\mathfrak p\) 处选一个不消失的 \(s\) 阶子式 \(d\)（\(s=0\) 时取 \(d=1\)），再局部化于 \(d\)。 可逆子块的行列消元把矩阵化为
	\[
	\begin{pmatrix}I_s&0\\0&B\end{pmatrix}.
	\]
	其构造是先把左上可逆块化为单位块，再减去其余行、列在这些主元上的分量。 在每个剩余域上总秩仍为 \(s\)，所以 \(B\) 的每个元素都落在每个素理想中。 由\cref{g5:prop:radical-prime-intersection}，这些元素是幂零元。 局部化保持约化性：若 \((a/t)^N=0\)，某个 \(u\) 使 \(ua^N=0\)，则 \((ua)^N=0\)，约化性给出 \(ua=0\)，故 \(a/t=0\)。 因而此处 \(B=0\)，得到 \(M_d\cong A_d^r\)。 每个点都有这样的邻域，故 \(M\) 局部自由。 双数环 \(k[\varepsilon]/(\varepsilon^2)\) 不约化；纤维中的消失只能说明剩余关系为幂零，不能迫使关系本身为零。
\end{solution}

\begin{exercise}\label{g5:ex:tensor-support-intersection}
	设 \(A\) 为交换含幺环，\(M,N\) 为有限生成 \(A\)-模。 证明
	\[
	\operatorname{Supp}_A(M\otimes_A N)
	=\operatorname{Supp}_A M\cap\operatorname{Supp}_A N.
	\]
\end{exercise}
\begin{solution}
	在每个素理想 \(\mathfrak p\) 处，局部化与张量积相容； 若 \(M_{\mathfrak p}\) 或 \(N_{\mathfrak p}\) 为零，张量积局部化也为零。 反过来，若二者都非零，Nakayama 引理保证 \(M\otimes_A\kappa(\mathfrak p)\) 和 \(N\otimes_A\kappa(\mathfrak p)\) 都是非零有限维向量空间。 它们的张量积维数为两个正维数之积，故非零。 又
	\[
	(M\otimes_A N)\otimes_A\kappa(\mathfrak p)
	\cong
	(M\otimes_A\kappa(\mathfrak p))
	\otimes_{\kappa(\mathfrak p)}
	(N\otimes_A\kappa(\mathfrak p)),
	\]
	所以 \((M\otimes_A N)_{\mathfrak p}\ne0\)。 两个包含合在一起得到等式。
\end{solution}
